% Chapter 5: Code generation using patterns

\chapter{Code generation using patterns}

\label{ch:fifth} % For referencing the chapter elsewhere, use \autoref{ch:name} 
\label{chapter:codeGeneration}

In the last two chapters we discussed how regular parallel pattern, referred to as \emph{algorithmic skeletons}, help to greatly simplify programming of modern parallel systems.
Parallel programming of multi-\GPU systems is considerably simplified without scarifying performance as shown by the evaluation in \autoref{chapter:skelcl-evaluation}.
In this chapter we address the second main challenge identified in \autoref{chapter:background}: \emph{Performance portability}.
We will present a novel approach for generating efficient and specialized code from compositions of algorithmic skeletons and similar high-level \emph{patterns}.
This systematic, rule-based approach provides performance portability across different types of modern parallel processors. 
In the following \autoref{chapter:codeGeneration-evaluation} we will show that this approach generates code matching the performance of highly tuned implementations on \CPUs and \GPUs using a set of example application.

We will start this chapter by looking at optimizations in \OpenCL and how applying them changes the source code of applications.
We will show how these optimizations can be hardware specific breaking portability and, furthermore, how optimizations for one particular hardware architecture can lead to poor performance on other hardware architectures.
All of this will motivate the necessity, when aiming for performance portability, for generating optimized and specialized code from a pattern-based high-level representation.
We will discuss the benefits an approach using code generation offers over a library based approach, like \SkelCL presented in \autoref{chapter:skelcl}.
Then we will give an overview of our approach and present it in more detail in the following sections.
\autoref{chapter:codeGeneration-evaluation} will present an evaluation of the approach using a set of application studies.

\section{A case study of OpenCL optimizations}
To understand the problems of performance and portability in the context of modern parallel processors we will study a simple application example: parallel reduction.
This discussed is based on the presentation \emph{``Optimizing Parallel Reduction in CUDA''} by \citeauthor{Harris2007}~\cite{Harris2007} where optimizations for implementing parallel reduction using \CUDA and targeting Nvidia \GPUs are presented.
Optimization guidelines like this exist from every hardware vendor giving application developers advice on how to exploit their hardware~\cite{CUDAProgrammingGuide,AMDProgrammingGuide,IntelGPUProgrammingGuide,IntelXeonProgrammingGuide}.

In \autoref{chapter:skelcl} we saw that we can express a parallel reduction using a single algorithmic skeleton: \reduce.
Here we look at how efficient \OpenCL implementation of this algorithm look like.
We are especially interested in gradually optimizing this application to see how beneficial the single optimization steps are.

We will first start by looking at the implementation and optimizations on one particular hardware architecture, using Nvidia's Fermi \GPU architecture as our example.
Then we will see how the optimized implementations perform on an AMD \GPU and Intel \CPU, to evaluate their portability.
Finally, we will use our observations to motivate the need for a pattern-based code generator for achieving performance portability.


\subsection{Optimizing parallel reduction for Nvidia's \GPU architectures}
For implementing the parallel reduction, or more precisely the parallel summation of an array, an approach using two \OpenCL kernels is used.
The first \OpenCL kernel is executed in parallel by multiple \OpenCL work groups, each producing a temporary result, then the second \OpenCL kernel is executed by a single \OpenCL work group producing the final result.
This strategy is applied as global synchronization, \ie, synchronization across work groups, is prohibited in \OpenCL.
By using the two kernel approach massive parallelism can be exploited in the first phase as multiple work groups operate concurrently on independent parts of the input array.
The synchronization inside the second work group ensures correctness when computing the final result.

We will follow the methodology in~\cite{Harris2007} and evaluate the performance of the different versions using the measured memory bandwidth as our metric.
This is reasonable as the reduction has a very low arithmetic intensity and its performance is, therefore, bound by the available memory bandwidth.

\paragraph{First OpenCL implementation}
\autoref{lst:reduce0} shows the first, basic version of the parallel reduction in \OpenCL.
This will be our starting point for the following optimizations.
In this implementation each work item first loads an element into the local memory (line~\ref{lst:reduce0:load}).
After a synchronization (line~\ref{lst:reduce0:firstBarrier}) all work items of a work group execute a for loop (lines~\ref{lst:reduce0:for:start}--\ref{lst:reduce0:for:end}) to perform a collective tree-based reduction.
In every iteration the if statement (line~\ref{lst:reduce0:if}) ensures that a declining number of work items remain active performing partial reductions in the shrinking reduction tree.
The second barrier in line~\ref{lst:reduce0:for:end} ensures that no race conditions occurs when accessing the shared local memory.
Finally, the work item in the work group with id zero writes back the computed result to the global memory in line~\ref{lst:reduce0:writeBack}.

The implementation presented in \autoref{lst:reduce0} is not straightforward to implement.
The application developer has to be familiar with the parallel execution model of \OpenCL to avoid race conditions and deadlocks.
For example, it is important that the second barrier in line~\ref{lst:reduce0:for:end} is placed \emph{after} and not \emph{inside} the if statement.
This is true, even though, work items not entering the if statement will never read from or write to memory and, therefore, can never be influenced by a race condition.
Nevertheless, \OpenCL requires all work item from a work group to execute all barrier statements in a kernel exactly the same number of times.
The application developer is responsible to ensure that this condition is met, otherwise a deadlock will occur.

Putting the programming difficulties aside this implementation does not provide high performance.
\citeauthor{Harris2007} reports that only 2.41\% of the theoretical peak memory bandwidth are utilized. 

\begin{lstlisting}[%                                                             
caption={First \OpenCL implementation of the parallel reduction achieving 2.41\% of the theoretical peak memory bandwith performance.},%
numbers=left,%
float=tb,
label={lst:reduce0}]
kernel
void reduce0(global float* g_idata, global float* g_odata,
             unsigned int n, local float* l_data) {
  unsigned int tid = get_local_id(0);
  unsigned int i   = get_global_id(0);
  l_data[tid] = (i < n) ? g_idata[i] : 0;$\label{lst:reduce0:load}$
  barrier(CLK_LOCAL_MEM_FENCE);$\label{lst:reduce0:firstBarrier}$

  // do reduction in local memory
  for(unsigned int s=1; s < get_local_size(0); s *= 2) {$\label{lst:reduce0:for:start}$
    if ((tid % (2*s)) == 0) {$\label{lst:reduce0:if}$
      l_data[tid] += l_data[tid + s]; }
    barrier(CLK_LOCAL_MEM_FENCE); }$\label{lst:reduce0:for:end}$

  // write result for this work group to global memory
  if (tid == 0) g_odata[get_group_id(0)] = l_data[0]; }$\label{lst:reduce0:writeBack}$
\end{lstlisting}

\FloatBarrier
\newpage

\paragraph{Avoid divergent branching}

\autoref{lst:reduce1} shows the second implementation.
The differences from the previous implementation are highlighted in the code.

When performing the collective tree-based reduction in a work group a shrinking number of work items remains active until the last remaining work item computes the result of the entire work group.
In the previous version the modulo operator was used to determine which work item remains active (see line~\ref{lst:reduce0:if} in \autoref{lst:reduce0}).
This leads to a situation were not consecutive work items remain active, but work items which id is divisible by 2, then by 4, then by 8, and so on.
In Nvidia's \GPU architectures 32 work items are grouped into a \emph{warps} and executed together, as described in \autoref{chapter:background}.
It is highly beneficial to program in a style where all 32 work items of a warp perform the same instructions, especially, to avoid divergent branching between work items of a warp.
Using the modulo operator to determine the active work items leads to highly divergent branching.
The second implementation in \autoref{lst:reduce1}, therefore, uses a different formula (line~\ref{lst:reduce1:index}) to determine the active work items which avoids divergent branching.

The performance improves by a factor of 2.33 as compared to the first implementation.
Still only 5.62\% of the theoretically available memory bandwidth are used by this version.

\begin{lstlisting}[%                                                             
caption={\OpenCL implementation of the parallel reduction avoiding divergent branching.
         This implementation achieves 5.62\% of the theoretical peak memory bandwidth performance.},%
numbers=left,%
float=tb,
label={lst:reduce1}]
kernel
void reduce1(global float* g_idata, global float* g_odata,
             unsigned int n, local float* l_data) {
  unsigned int tid = get_local_id(0);
  unsigned int i   = get_global_id(0);
  l_data[tid] = (i < n) ? g_idata[i] : 0;
  barrier(CLK_LOCAL_MEM_FENCE);

  for(unsigned int s=1; s < get_local_size(0); s *= 2) {
      // continuous work items remain active
      $\strut$@int index = 2 * s * tid;@$\label{lst:reduce1:index}$
      if (@index < get_local_size(0)@) {$\label{lst:reduce1:if}$
          l_data[index] += l_data[index + s]; }$\label{lst:reduce1:read}$
      barrier(CLK_LOCAL_MEM_FENCE); }

  if (tid == 0) g_odata[get_group_id(0)] = l_data[0]; }
\end{lstlisting}


\paragraph{Avoid interleaved addressing}

\autoref{lst:reduce2} shows the third implementation.
The differences from the previous implementation are highlighted in the code.

On modern \GPUs the fast local memory is organized in multiple \emph{banks}.
When two, or more, work items simultaneously access memory locations in the same bank a \emph{bank conflict} occurs which means that all memory requests are answered sequentially and not in parallel as usual.
This is described in more detail in~\autoref{chapter:background}.
The previous two implementation use an access pattern for the local memory which makes bank conflicts likely.
When reading memory from the local memory in line~\ref{lst:reduce1:read} of \autoref{lst:reduce1} each work item reads interleaved from two locations: \lstinline!index! and \lstinline!index+s!.
The third implementation in \autoref{lst:reduce2} avoids this problematic local memory access pattern.
Instead a sequential access pattern is used where consecutive work items access consecutive memory locations.
This is achieved by directly using the local id as index together with a different definition of \lstinline!s! in line~\ref{lst:reduce2:s}.
This better access pattern requires the reduction operation to be commutative, as the order of element is not respected when reading from global memory.

The performance improves by a factor of 2.01 as compared to the previous implementation and 4.68 to the initial implementation.
With this version 11.27\% of the theoretically available memory bandwidth are used.


\begin{lstlisting}[%                                                             
caption={\OpenCL implementation of the parallel reduction avoiding local memory bank conflicts.
         This implementation achieves 11.27\% of the theoretical peak memory bandwidth performance.},%
numbers=left,%
float=tb,
label={lst:reduce2}]
kernel
void reduce2(global float* g_idata, global float* g_odata,
             unsigned int n, local float* l_data) {
  unsigned int tid = get_local_id(0);
  unsigned int i   = get_global_id(0);
  l_data[tid] = (i < n) ? g_idata[i] : 0;
  barrier(CLK_LOCAL_MEM_FENCE);

  // process elements in different order
  // requires commutativity!
  for(@unsigned int s=get_local_size(0)/2; s>0; s>>=1@) {$\label{lst:reduce2:s}$
      if (@tid < s@) {
          l_data[@tid@] += l_data[@tid@ + s]; }
      barrier(CLK_LOCAL_MEM_FENCE); }

  if (tid == 0) g_odata[get_group_id(0)] = l_data[0]; }
\end{lstlisting}

\paragraph{Increase computational intensity per work item}

\autoref{lst:reduce3} shows the fourth implementation.
The differences from the previous implementation are highlighted in the code.

For the previous versions each work item loads one element from global into local memory before the work group collectively performs a tree-based reduction.
That means that half of the work items are idle after performing a single copy operation, which is highly wasteful.
The fourth implementation in \autoref{lst:reduce3} avoids this by letting each work item load two elements from global memory, perform an addition, and store the computed result in local memory.
Assuming the same input size this reduces the number of work items to start by half and, therefore, increases the computational intensity for every work item.

The performance improves by a factor of 1.78 as compared to the previous implementation and 8.34 to the initial implementation.
With this version 20.11\% of the theoretically available memory bandwidth are used.


\begin{lstlisting}[%                                                             
caption={\OpenCL implementation of the parallel reduction. Each work item performs an addition when loading data from global memory.
         This implementation achieves 20.11\% of the theoretical peak memory bandwidth performance.},%
numbers=left,%
float=tb,
escapechar=\`,
label={lst:reduce3}]
kernel
void reduce3(global float* g_idata, global float* g_odata,
             unsigned int n, local float* l_data) {
  unsigned int tid = get_local_id(0);
  $\strut$@unsigned int i = get_group_id(0) * (get_local_size(0)*2)@
                                   $\strut$@+ get_local_id(0);@
  l_data[tid] = (i < n) ? g_idata[i] : 0;
  // performs first addition during loading
  $\strut$@if (i + get_local_size(0) < n)@
  $\strut$@    l_data[tid] += g_idata[i+get_local_size(0)];@
  barrier(CLK_LOCAL_MEM_FENCE);

  for(unsigned int s=get_local_size(0)/2; s>0; s>>=1) {
      if (tid < s) {
          l_data[tid] += l_data[tid + s]; }
      barrier(CLK_LOCAL_MEM_FENCE); }

  if (tid == 0) g_odata[get_group_id(0)] = l_data[0]; }
\end{lstlisting}

\newpage

\paragraph{Avoid synchronization inside a warp}

\autoref{lst:reduce4} shows the fifth implementation.
The differences from the previous implementation are highlighted in the code.

Wraps are the fundamental execution unit in Nvidia's \GPU architectures, as explained in \autoref{chapter:background}.
32 work items are grouped together to form a warp based on their id, \ie, work items with id 0 -- 31 are grouped into the first warp, work items with id 32 -- 63 into the second warp, \etc.
All work items grouped in a warp are guaranteed to be executed together in a lock-step manner:
all work items in the same warp execute the same instruction simultaneously.
Because of this hardware behaviour, no barrier synchronization is required between instructions inside a single warp.
The fifth implementation in \autoref{lst:reduce4} takes advantage of this.
The for loop performing the tree-based reduction is left early at the stage when only 32 work items remain active (see line~\ref{lst:reduce4:for}).
Lines~\ref{lst:reduce4:warp:start}--\ref{lst:reduce4:warp:end} shows the special code which performs the rest of the tree-base reduction without any barrier synchronization.
The code shown here effectively unrolled the last six iterations of the for loop in line~\ref{lst:reduce4:for}.
As warps are specific to Nvidias \GPU architectures, this implementation is not portable \OpenCL code and will produce incorrect results on other \OpenCL devices.

The performance improves by a factor of 1.8 as compared to the previous implementation and 15.01 to the initial implementation.
With this version 36.21\% of the theoretically available memory bandwidth are used.

\begin{lstlisting}[%                                                             
caption={\OpenCL implementation of the parallel reduction.
         Synchronization inside a warp is avoided by unrolling the loop for the last 32 work items.
         This implementation achieves 36.21\% of the theoretical peak memory bandwidth performance.},%
numbers=left,%
float=p,
label={lst:reduce4}]
kernel
void reduce4(global float* g_idata, global float* g_odata,
             unsigned int n,local volatile float* l_data){
  unsigned int tid = get_local_id(0);
  unsigned int i = get_group_id(0) * (get_local_size(0)*2)
                                   + get_local_id(0);
  l_data[tid] = (i < n) ? g_idata[i] : 0;
  if (i + get_local_size(0) < n) 
      l_data[tid] += g_idata[i+get_local_size(0)];  
  barrier(CLK_LOCAL_MEM_FENCE);

  // prevent further unrolling (see next version)
  $\strut$@#pragma unroll 1@$\label{lst:reduce4:pragma}$
  for(unsigned int s=get_local_size(0)/2; s>@32@; s>>=1) {$\label{lst:reduce4:for}$
      if (tid < s) {
          l_data[tid] += l_data[tid + s]; }
      barrier(CLK_LOCAL_MEM_FENCE); }

  // unroll for last 32 active work items
  // no synchronization required on NVIDIA GPUs
  // this is not protable OpenCL code!
  $\strut$@if (tid < 32) {@$\label{lst:reduce4:warp:start}$
  $\strut$@  if (WG_SIZE >= 64) { l_data[tid] += l_data[tid+32]; }@
  $\strut$@  if (WG_SIZE >= 32) { l_data[tid] += l_data[tid+16]; }@
  $\strut$@  if (WG_SIZE >= 16) { l_data[tid] += l_data[tid+ 8]; }@
  $\strut$@  if (WG_SIZE >=  8) { l_data[tid] += l_data[tid+ 4]; }@
  $\strut$@  if (WG_SIZE >=  4) { l_data[tid] += l_data[tid+ 2]; }@
  $\strut$@  if (WG_SIZE >=  2) { l_data[tid] += l_data[tid+ 1]; }@
  $\strut$@}@$\label{lst:reduce4:warp:end}$

  if (tid == 0) g_odata[get_group_id(0)] = l_data[0]; }
\end{lstlisting}

\paragraph{Complete loop unrolling}

\autoref{lst:reduce5} shows the sixth implementation.
The differences from the previous implementation are highlighted in the code.

In the previous implementation we made a special case for the last six iterations of the for loop and provided special code handling each iteration separately.
This is a general optimization strategy known as \emph{loop unrolling}.
Loop unrolling can be beneficial because variables and branches required by the loop can be avoided.
In the sixth implementation shown in \autoref{lst:reduce5} the for loop has been removed entirely and the replaced by three if statement (line~\ref{lst:reduce5:if:1},~\ref{lst:reduce5:if:2}, and~\ref{lst:reduce5:if:3}).
Each if statement replaces one iteration of the loop.
This code assumes that \code{WG\_SIZE} is a compile time constant and, therefore, the if statements will be evaluated at compile time, avoiding costly branches at runtime.
Different to the previous optimization we still have to provide a barrier to ensure correct synchronization, as multiple warps are involved here.

The performance improves by a factor of 1.41 as compared to the previous implementation and 21.16 to the initial implementation.
With this version 56.92\% of the theoretically available memory bandwidth are used.

\begin{lstlisting}[%                                                             
caption={\OpenCL implementation of the parallel reduction with a completly unrolled loop.
         This implementation achieves 56.92\% of the theoretical peak memory bandwidth performance.},%
numbers=left,%
float=p,
label={lst:reduce5}]
kernel
void reduce5(global float* g_idata, global float* g_odata,
             unsigned int n,local volatile float* l_data){
  unsigned int tid = get_local_id(0);
  unsigned int i = get_group_id(0) * (get_local_size(0)*2)
                                   + get_local_id(0);
  l_data[tid] = (i < n) ? g_idata[i] : 0;
  if (i + get_local_size(0) < n) 
      l_data[tid] += g_idata[i+get_local_size(0)];  
  barrier(CLK_LOCAL_MEM_FENCE);

  // unroll for loop entirely
  $\strut$@if (WG_SIZE >= 512) {@$\label{lst:reduce5:if:1}$
  $\strut$@    if (tid < 256) { l_data[tid] += l_data[tid+256]; }@
  $\strut$@    barrier(CLK_LOCAL_MEM_FENCE); }@
  $\strut$@if (WG_SIZE >= 256) {@$\label{lst:reduce5:if:2}$
  $\strut$@    if (tid < 128) { l_data[tid] += l_data[tid+128]; }@
  $\strut$@    barrier(CLK_LOCAL_MEM_FENCE); }@
  $\strut$@if (WG_SIZE >= 128) {@$\label{lst:reduce5:if:3}$
  $\strut$@    if (tid <  64) { l_data[tid] += l_data[tid+ 64]; }@
  $\strut$@    barrier(CLK_LOCAL_MEM_FENCE); }@
  
  if (tid < 32) {
    if (WG_SIZE >= 64) { l_data[tid] += l_data[tid+32]; }
    if (WG_SIZE >= 32) { l_data[tid] += l_data[tid+16]; }
    if (WG_SIZE >= 16) { l_data[tid] += l_data[tid+ 8]; }
    if (WG_SIZE >=  8) { l_data[tid] += l_data[tid+ 4]; }
    if (WG_SIZE >=  4) { l_data[tid] += l_data[tid+ 2]; }
    if (WG_SIZE >=  2) { l_data[tid] += l_data[tid+ 1]; } }
  
  if (tid == 0) g_odata[get_group_id(0)] = l_data[0]; }
\end{lstlisting}


\paragraph{Fully optimized implementation}

\autoref{lst:reduce6} shows the last and fully optimized implementation.
The differences from the previous implementation are highlighted in the code.

One of the optimizations applied earlier was to increase the computational intensity for each single work item by performing two loads and an addition instead of a single load.
This final version applies the same idea, but performing multiple additions per work item before the collective tree-based reduction with the entire work group.
This has indeed two advantages:
first, the algorithmic intensity is increased, \ie, each work item is doing more work,
and, second, performing the summation sequentially by a single work item does not require costly synchronizations.
The fully optimized implementation is shown in \autoref{lst:reduce6} with the changes highlighted.
A while loop has been introduced (see line~\ref{lst:reduce6:while}) which, in every iteration, loads two elements from global memory and adds them to the local memory.
No synchronization is required here as each work item operates independently on different memory locations.
The way the memory is accessed ensures that memory accesses will be coalesced when reading from global memory.

The performance improves by a factor of 1.42 as compared to the previous implementation and 30.04 to the initial implementation.
With this version 72.53\% of the theoretically available memory bandwidth are used.

\begin{lstlisting}[%                                                             
caption={Fully optimized \OpenCL implementation of the parallel reduction.
         This implementation achieves 72.53\% of the theoretical peak memory bandwidth performance.},%
numbers=left,%
float=p,
label={lst:reduce6}]
kernel
void reduce6(global float* g_idata, global float* g_odata,
             unsigned int n,local volatile float* l_data){
  unsigned int tid = get_local_id(0);
  unsigned int i = get_group_id(0) * (get_local_size(0)*2)
                                   + get_local_id(0);
  $\strut$@unsigned int gridSize = WG_SIZE*2*get_num_groups(0);@
  $\strut$@l_data[tid] = 0;@

  // multiple elements are reduced per work item
  $\strut$@while (i < n) { l_data[tid] += g_idata[i];@$\label{lst:reduce6:while}$
  $\strut$@                if (i + WG_SIZE < n)@
  $\strut$@                  l_data[tid] += g_idata[i+WG_SIZE];@
  $\strut$@                i += gridSize; }@
  barrier(CLK_LOCAL_MEM_FENCE);

  if (WG_SIZE >= 512) {
      if (tid < 256) { l_data[tid] += l_data[tid+256]; }
      barrier(CLK_LOCAL_MEM_FENCE); }
  if (WG_SIZE >= 256) {
      if (tid < 128) { l_data[tid] += l_data[tid+128]; }
      barrier(CLK_LOCAL_MEM_FENCE); }
  if (WG_SIZE >= 128) {
      if (tid <  64) { l_data[tid] += l_data[tid+ 64]; }
      barrier(CLK_LOCAL_MEM_FENCE); }
  
  if (tid < 32) {
    if (WG_SIZE >= 64) { l_data[tid] += l_data[tid+32]; }
    if (WG_SIZE >= 32) { l_data[tid] += l_data[tid+16]; }
    if (WG_SIZE >= 16) { l_data[tid] += l_data[tid+ 8]; }
    if (WG_SIZE >=  8) { l_data[tid] += l_data[tid+ 4]; }
    if (WG_SIZE >=  4) { l_data[tid] += l_data[tid+ 2]; }
    if (WG_SIZE >=  2) { l_data[tid] += l_data[tid+ 1]; } }
  
  if (tid == 0) g_odata[get_group_id(0)] = l_data[0]; }
\end{lstlisting}


\paragraph{Conclusions}
We can draw valuable conclusions from studying these \OpenCL source codes.

The first conclusion is, that implementing these optimizations is not intuitive and straightforward.
It requires experience as well as knowledge and reasoning about the target hardware architecture, in this case Nvidia's Fermi \GPU architecture:
\autoref{lst:reduce1} requires the understanding of the problem of \emph{branch divergence}, \autoref{lst:reduce2} requires knowledge about the organization of local memory and \emph{bank conflicts}, \autoref{lst:reduce3} and \autoref{lst:reduce6} require reasoning about the \emph{computational intensity} of work items, \autoref{lst:reduce4} requires understanding of \emph{warps} and their execution by the hardware, \autoref{lst:reduce5} requires experience with \emph{loop unrolling} techniques, and, finally, \autoref{lst:reduce6} required knowledge about the organization of global memory and \emph{memory coalescing}.
The author wants to empathize that these are \emph{additional} burdens for the application developer on top of implementing a functionally correct version facing widely recognized problems of parallel programming like dealing with race conditions and deadlocks.

The second conclusion we can draw is, that performing these optimizations on modern parallel architectures is highly beneficial.
The first unoptimized version did only utilize about 2.5\% of the available memory bandwidth, while the fully optimized version utilizes a more reasonable 72.53\%.
Applying all optimizations improved the performance by a factor of $30$ while utilizing the exact same hardware.
Modern parallel processors are often chosen as target architecture because of their high theoretical performance.
Turning the theoretical performance into practical performance by applying these optimizations is, therefore, \emph{essential} for most users.





\subsection{Portability of the optimized parallel reduction}
After we have established how crucial but hard to achieve optimizations are, we now will investigate their portability.
To do so, we run the code shown in \autoref{lst:reduce0}--\ref{lst:reduce6} on three different hardware architectures:
Nvidia's Fermi \GPU architecture~\cite{} which we have used in our analysis in the last section, AMD's Tahiti \GPU architecture~\cite{}, and Intel's Nehalem \CPU architecture~\cite{}.
\autoref{fig:reduce:overview} shows the performance numbers separated for each architecture.
We use, as before, the memory bandwidth as our performance metric and show percentages of the theoretical memory bandwidth of each respective hardware architecture.
As a practical comparison we also show performance numbers for the parallel reduction using highly tuned, architecture specific implementations of the BLAS library.
We use the cuBLAS library~\cite{} for the Nvidia \GPU, the clBLAS library~\cite{} for the AMD \GPU, and the MKL library~\cite{} for the Intel \CPU.
Each library implementation is provided by the corresponding hardware vendor.

The results in \autoref{fig:reduce:overview} show that the optimizations discussed in the previous section are not portable. 
On each architecture a different version performs best.
Some versions actually produce incorrect results on the \CPU due to the warp specific optimization.
This optimization happens to be valid on AMD's \GPU architecture as well, as there exists a similar concept to warps called \emph{wavefront}~\cite{}.





\subsection{The need for a pattern-based code generator}

Our goal in this chapter is to achieve \emph{performance portability}, \ie, to achieve high performance for a given application across a set of different parallel processors.
Achieving performance portability with traditional approaches is hard, as there is a tension between achieving the highest performance possible and code portability and maintainability.
Traditionally, \eg, in C or \OpenCL, programmers tune their implementations towards a particular hardware using hardware-specific optimizations to achieve the highest performance possible.
This reduces portability, maintainability, and clarity of the code:
multiple versions have to be maintained and non-obvious optimizations make the code hard to understand and to reason about.

We argue that parallel pattern can overcome this fundamental conflict as they declaratively specify the desired algorithmic behavior rather than encode a particular implementation which might offer suboptimal performance on some hardware architectures.
The parallel pattern can be implemented in different ways optimized towards particular hardware architectures.
If the underlying hardware for an application implemented with parallel patterns is changed, the most optimized implementation for the new hardware can be chosen.

\todo{...}
- This approach imposes challenges

1) optimized implementation of every pattern on every new hardware

2) composition and nesting is difficult to handle.

3) optimal implementation might be application and context specific

- examples, simple map, reduce also, in dot prod, in matrix mult

- Hard or impossible to overcome with library approach.

% what is our approach
The root of the problem lies in a gap in the system stack between high-level algorithmic concepts on the one hand and low-level hardware paradigms on the other hand.
Previous work has proposed ad-hoc solutions to target specific hardware architectures.
We propose to bridge this gap by defining a set of rewrite rules which systematically translates high-level algorithmic concepts into low-level hardware paradigms, both expressed as functional patterns.
The rewrite rules are used to systematically derive semantically equivalent low-level expressions from high-level algorithm expressions written by the application developer.
Once derived, we can automatically generate high performance code based on these expressions.
The next section introduces on overview of our approach.
 


\section{Overview}

gap in system stack

\from{PACT begin}
\paragraph{Motivation PACT}
The overview of our approach is presented in Figure~\ref{fig:highlevel}.
The programmer expresses his problem by writing a \emph{high-level expression} composed of \emph{algorithmic patterns}.
Using a rewrite rule system, we automatically lower this high-level expression into a \emph{low-level expression} consisting of \emph{hardware patterns}.
This rewrite stage explores the algorithmic and optimization choices in the high-level expression.
The generated low-level expression is then fed into our code generator that emits an \emph{OpenCL program}.
This program is finally compiled to machine code by the vendor provided OpenCL compiler.

The biggest advantage of our approach is that there is no analysis or optimizations performed in the code generator.
All optimization decisions are automatically made earlier on in the rule rewriting system.
This results in a clear separation of concern between the high-level patterns used by the programmer and the low-level hardware paradigms that enable performance portability.

\begin{figure}[t]
\centering
\includegraphics[width=\linewidth]{overviewPatternCodeGeneration}
\vspace{-15pt}
\caption{
The programmer expresses his problem with high-level algorithmic patterns.
A rule rewriting system explores the space of implementations and lower the high-level expression.
Finally, the code generator produces OpenCL code by mapping the low-level hardware patterns directly to the OpenCL programming model representing the hardware paradigms.}
\label{fig:highlevel}
\end{figure}


\subsection{Example}

\begin{figure}[t]
\centering

\begin{subfigure}[b]{.85\linewidth}
\begin{lstlisting}[mathescape,numbers=left]
def mul3(x) = x * 3    // user-defined function
def vectorMul = map(mul3)        // map pattern
\end{lstlisting}
\caption{\textbf{High-level expression} written by the programmer.}
\label{fig:codeex:map}
\end{subfigure}

\vspace{-5pt}
\begin{minipage}{0.1\linewidth}
\vspace{0pt}
\centering
\begin{tikzpicture}[ultra thick]
 \draw [black,   -latex      ] (0,0.5) -- (0,0) node [] {};
\end{tikzpicture}
\end{minipage}
\begin{minipage}{0.25\linewidth}
\vspace{-5pt}
\centering
\textbf{rewrite rules}
\end{minipage}
\begin{minipage}{0.1\linewidth}
\vspace{0pt}
\centering
\begin{tikzpicture}[ultra thick]
 \draw [black,   -latex      ] (0,0.5) -- (0,0) node [] {};
\end{tikzpicture}
\end{minipage}

\begin{subfigure}[b]{\linewidth}
\centering
\begin{minipage}{.85\linewidth}
\begin{lstlisting}[mathescape,numbers=left]
def mul3(x) = x * 3
def vectorMul = join $\circ$ map-workgroup(
                  join $\circ$ map-local(
                    vectorize-4(mul3)
                  ) $\circ$ split-4
                ) $\circ$ split-1024
\end{lstlisting}
\end{minipage}
\caption{\textbf{Low-level expression} automatically derived using rewrite rules.}
\label{fig:codeex:impl}
\end{subfigure}

\vspace{-5pt}
\begin{minipage}{0.1\linewidth}
\vspace{0pt}
\centering
\begin{tikzpicture}[ultra thick]
 \draw [black,   -latex      ] (0,0.5) -- (0,0) node [] {};
\end{tikzpicture}
\end{minipage}
\begin{minipage}{0.26\linewidth}
\vspace{-5pt}
\centering
\textbf{code generator}
\end{minipage}
\begin{minipage}{0.1\linewidth}
\vspace{0pt}
\centering
\begin{tikzpicture}[ultra thick]
 \draw [black,   -latex      ] (0,0.5) -- (0,0) node [] {};
\end{tikzpicture}
\end{minipage}

\begin{subfigure}[b]{\linewidth}
\centering
\begin{minipage}{.85\textwidth}
\begin{lstlisting}[mathescape,numbers=left]
int4 mul3(int4 x) { return x * 3; }
kernel map_mul3(global int* in,out, int len) {
  // division into workgroup by chuncks of 1024
  for (int i=get_group_id; i < len/1024;
       i+=get_num_groups) {
    global int* grp_in  = in+(i*1024);
    global int* grp_out = out+(i*1024);
    // division into threads by chunks of 4
    for (int j=get_local_id; j < 1024/4;
         j+=get_local_size) {
      global int* lcl_in  = grp_in+(j*4);
      global int* lcl_out = grp_out+(j*4);
      // vectorization with vector width of 4
      global int4* in_vec4 = (int4*) lcl_in;
      global int4* out_vec4 = (int4*) lcl_out;
      *out_vec4 = mul3(*in_vec4);      
} } }  
\end{lstlisting}
\end{minipage}
\caption{\textbf{OpenCL program} automatically produced by our code generator.}
\label{fig:codeex:ocl}
\end{subfigure}

\caption{Pseudo-code representing a vector scaling program.
The user simply maps the \texttt{mul3} function over the elements of the input array~(\subref{fig:codeex:map}).
This high-level expression is automatically transformed into a low-level expression~(\subref{fig:codeex:impl}) using rewrite rules.
Finally, our code generator turns the low-level expression into an OpenCL program~(\subref{fig:codeex:ocl}).}
\label{fig:codeex}
\end{figure}

We now present the advantages of our approach with a simple vector scaling example shown in Figure~\ref{fig:codeex}.
The user expresses the computation by writing a high-level expression using our \pat{map} algorithmic pattern as shown in Figure~\ref{fig:codeex:map}.
This coding style is similar to functional and dataflow programming.

Our technique first rewrites automatically the user provided high-level expression into something closer to the OpenCL programming model.
This is achieved by applying the rewrite rules presented later in Section~\ref{rules}.
Figure~\ref{fig:codeex:impl} shows one possible derivation of the original high-level map expression where the $\circ$ operator represents function composition, \ie $f \circ g(x) = f(g(x))$.
Starting from the last line, we first split the input into chunks of 1024 elements.
Each chunk is mapped onto a workgroup with the \pat{map-workgroup} low-level pattern (line~2).
Within a workgroup (lines~3--5), we further split the elements into chunks of 4 (line~5), each mapped to a local thread via the \pat{map-local} low-level pattern (line~3).
Each local thread is responsible for processing 4 elements.
Finally, the \pat{vectorize-4} pattern (line~4) implies that the user defined function \texttt{mul3} is vectorized.
The exact meaning of our patterns will be given later in Section~\ref{pattern}.

The last step consists of traversing the low-level expression and generating OpenCL code (see Figure~\ref{fig:codeex:ocl}) for each low-level pattern encountered.
The two map patterns generate the for loops (line~4--5 and~9--10) that iterate over the input array assigning work to the workgroups and local threads.
The information of how many chunks each workgroup and thread processes comes from the corresponding \pat{split}.
In line~16 the vectorized version of the user defined \texttt{mul3} function (defined in line~1) is finally applied to the input array.

To summarize, our approach is able to automatically generate OpenCL code starting from a high-level representation of a program.
This is achieved by automatically lowering the high-level expression into a low-level form suitable for code generation.
The next two sections present our high-level and low-level patterns, the code generation mechanism and the rewrite rules in more details.
\from{PACT end}

\from{PPoPP begin}
\paragraph{Motivation PPoPP}
The overview of our approach is presented in Figure~\ref{fig:highlevel}.
The programmer writes a \emph{high-level expression} composed of \emph{algorithmic patterns}.
Using a rewrite rule system, we systematically lower this high-level expression into a \emph{low-level expression} consisting of \emph{OpenCL patterns}.
In this rewrite stage algorithmic and optimization choices in the high-level expression can be explored.
The generated low-level expression is then fed into our code generator that emits an \emph{OpenCL program}.
This program is finally compiled to machine code by the vendor provided OpenCL compiler.

% One big advantage of our approach is that there is no analysis or optimizations performed in the code generator.
% All optimization decisions are automatically made earlier on in the rule rewriting system.
% This results in a clear separation of concern between the high-level patterns used by the programmer and the low-level hardware paradigms that enable performance portability.

\begin{figure}[t]
\centering
\includegraphics[width=\linewidth]{overviewPatternCodeGeneration}
\vspace{-19pt}
\caption{
Overview of our system.
The programmer expresses the problem with high-level algorithmic patterns.
These are systematically transformed into low-level OpenCL patterns using a rule rewriting system.
OpenCL code is generated by mapping the low-level patterns directly to the OpenCL programming model representing hardware paradigms.
\vspace{-1em}}
\label{fig:highlevel}
\end{figure}





% \subsection{Example}

\begin{figure}[t]
\centering

\begin{subfigure}[b]{.85\linewidth}%{0.7\linewidth}
\begin{lstlisting}[mathescape,numbers=left]
def mul3(x) = x * 3    // user-defined function
def vectorScal = map(mul3)       // map pattern
\end{lstlisting}
\caption{\textbf{High-level expression} written by the programmer.}
\label{fig:codeex:map}
\end{subfigure}

\vspace{-10pt}
\begin{minipage}{0.1\linewidth}
\vspace{0pt}
\centering
\begin{tikzpicture}[ultra thick]
 \draw [black,   -latex      ] (0,0.5) -- (0,0) node [] {};
\end{tikzpicture}
\end{minipage}
\begin{minipage}{0.25\linewidth}
\vspace{-5pt}
\centering
\textbf{rewrite rules}
\end{minipage}
\begin{minipage}{0.1\linewidth}
\vspace{0pt}
\centering
\begin{tikzpicture}[ultra thick]
 \draw [black,   -latex      ] (0,0.5) -- (0,0) node [] {};
\end{tikzpicture}
\end{minipage}

\vspace{0pt}
\begin{subfigure}[b]{\linewidth}
\centering
\begin{minipage}{.85\linewidth}%{0.85\textwidth}
\begin{lstlisting}[mathescape,numbers=left]
def mul3(x) = x * 3
def vectorScal = join $\circ$ map-workgroup(
                  asScalar $\circ$ map-local(
                    vectorize-4(mul3)
                  ) $\circ$ asVector-4
                 ) $\circ$ split-1024
\end{lstlisting}
\end{minipage}
\caption{\textbf{Low-level expression} systematically derived using rewrite rules.}
\label{fig:codeex:impl}
\end{subfigure}

\vspace{-10pt}
\begin{minipage}{0.1\linewidth}
\vspace{0pt}
\centering
\begin{tikzpicture}[ultra thick]
 \draw [black,   -latex      ] (0,0.5) -- (0,0) node [] {};
\end{tikzpicture}
\end{minipage}
\begin{minipage}{0.26\linewidth}
\vspace{-5pt}
\centering
\textbf{code generator}
\end{minipage}
\begin{minipage}{0.1\linewidth}
\vspace{0pt}
\centering
\begin{tikzpicture}[ultra thick]
 \draw [black,   -latex      ] (0,0.5) -- (0,0) node [] {};
\end{tikzpicture}
\end{minipage}

\vspace{0pt}
\begin{subfigure}[b]{\linewidth}
\centering
\begin{minipage}{.85\textwidth}
\begin{lstlisting}[mathescape,numbers=left]
int4 mul3(int4 x) { return x * 3; }
kernel vectorScal(global int* in,out, int len){
 for (int i=get_group_id; i < len/1024;
      i+=get_num_groups) {
  global int* grp_in  = in+(i*1024);
  global int* grp_out = out+(i*1024);
  for (int j=get_local_id; j < 1024/4;
       j+=get_local_size) {
    global int4* in_vec4 =(int4*)grp_in+(j*4);
    global int4* out_vec4=(int4*)grp_out+(j*4);
    *out_vec4 = mul3(*in_vec4);      
} } }  
\end{lstlisting}
\end{minipage}
\caption{\textbf{OpenCL program} produced by our code generator.}
\label{fig:codeex:ocl}
\end{subfigure}
\vspace{-20pt}
\caption{
Pseudo-code representing vector scaling.
The user simply maps the \code{mul3} function over the elements of the input array~(\subref{fig:codeex:map}).
This high-level expression is systematically transformed into a low-level expression~(\subref{fig:codeex:impl}) using rewrite rules.
Finally, our code generator turns the low-level expression into an OpenCL program~(\subref{fig:codeex:ocl}).
}
  \label{fig:codeex}
\end{figure}

We now illustrate the advantages of our approach using a simple vector scaling example shown in Figure~\ref{fig:codeex}.
The user expresses the computation by writing a high-level expression using our \pat{map} algorithmic pattern as shown in Figure~\ref{fig:codeex:map}.
This coding style is similar to functional and dataflow programming.

Our technique first rewrites the user provided high-level expression into something closer to the OpenCL programming model.
This is achieved by applying the rewrite rules presented later in Section~\ref{rules}.
Figure~\ref{fig:codeex:impl} shows one possible derivation of the original high-level expression where the $\circ$ operator represents function composition, \ie $f \circ g(x) = f(g(x))$.
Starting from the last line, we first split the input into chunks of 1024 elements.
Each chunk is mapped onto a group of threads, called \emph{workgroup}, with the \pat{map-workgroup} low-level pattern (line~2).
Within a workgroup (lines~3--5), we vectorize the elements (line~5), each mapped to a local thread inside a workgroup via the \pat{map-local} low-level pattern (line~3).
Each local thread now processes 4 elements, enclosed in a vector type.
Finally, the \pat{vectorize-4} pattern (line~4) implies that the user defined function \code{mul3} is vectorized.
The exact meaning of our patterns will be given later in Section~\ref{pattern}.

The last step consists of traversing the low-level expression and generating OpenCL code for each low-level pattern encountered (Figure~\ref{fig:codeex:ocl}).
The two map patterns generate the for loops (line~3--4 and~7--8) that iterate over the input array assigning work to the workgroups and local threads.
The information of how many chunks each workgroup and thread processes comes from the corresponding \pat{split}.
In line~11 the vectorized version of the user defined \code{mul3} function (defined in line~1) is finally applied to the input array.

To summarize, our approach is able to generate OpenCL code starting from a high-level representation of a program.
This is achieved by systematically lowering the high-level expression into a low-level form suitable for code generation.
The next two sections present our high-level and low-level patterns, the code generation mechanism and the rewrite rules in more details.
\from{PPoPP end}

\from{PLDI begin}
\paragraph{Motivation PLDI}
The overview of our approach is presented in Figure~\ref{fig:highlevel}.
The programmer writes a \emph{high-level expression} composed of \emph{algorithmic patterns}.
Using a rewrite rule system, we lower this high-level expression into a \emph{low-level expression} consisting of \emph{OpenCL patterns}.
In this rewrite stage algorithmic and optimization choices in the high-level expression can be explored.
The generated low-level expression is then fed into our code generator that emits an \emph{OpenCL program} compiled to machine code by the vendor provided OpenCL compiler.

\begin{figure}[t]
\centering
\includegraphics[width=\linewidth]{overviewPatternCodeGeneration}
\vspace{-19pt}
\caption{
Overview of our system.
The programmer expresses the problem with high-level algorithmic patterns.
These are systematically transformed into low-level patterns using a rule rewriting system.
OpenCL code is generated by mapping the low-level patterns directly to the OpenCL programming model representing hardware paradigms.
\vspace{-1em}}
\label{fig:highlevel}
\end{figure}





% \subsection{Example}

\begin{figure}[t]
\centering

\begin{subfigure}[b]{.85\linewidth}%{0.7\linewidth}
\begin{lstlisting}[mathescape,numbers=left]
def mul3(x) = x * 3    // user-defined function
def vectorScal = map(mul3)       // map pattern
\end{lstlisting}
\caption{\textbf{High-level expression} written by the programmer.}
\label{fig:codeex:map}
\end{subfigure}

\vspace{-10pt}
\begin{minipage}{0.1\linewidth}
\vspace{0pt}
\centering
\begin{tikzpicture}[ultra thick]
 \draw [black,   -latex      ] (0,0.5) -- (0,0) node [] {};
\end{tikzpicture}
\end{minipage}
\begin{minipage}{0.25\linewidth}
\vspace{-5pt}
\centering
\textbf{rewrite rules}
\end{minipage}
\begin{minipage}{0.1\linewidth}
\vspace{0pt}
\centering
\begin{tikzpicture}[ultra thick]
 \draw [black,   -latex      ] (0,0.5) -- (0,0) node [] {};
\end{tikzpicture}
\end{minipage}

\vspace{-5pt}
\begin{subfigure}[b]{\linewidth}
\centering
\begin{minipage}{.85\linewidth}%{0.85\textwidth}
\begin{lstlisting}[mathescape,numbers=left]
def mul3(x) = x * 3
def vectorScal = join $\circ$ map-workgroup(
                  asScalar $\circ$ map-local(
                    vect-4(mul3)
                  ) $\circ$ asVector-4
                 ) $\circ$ split-1024
\end{lstlisting}
\end{minipage}
\caption{\textbf{Low-level expression} derived using rewrite rules.}
\label{fig:codeex:impl}
\end{subfigure}

\vspace{-5pt}
\begin{minipage}{0.1\linewidth}
\vspace{0pt}
\centering
\begin{tikzpicture}[ultra thick]
 \draw [black,   -latex      ] (0,0.5) -- (0,0) node [] {};
\end{tikzpicture}
\end{minipage}
\begin{minipage}{0.26\linewidth}
\vspace{-5pt}
\centering
\textbf{code generator}
\end{minipage}
\begin{minipage}{0.1\linewidth}
\vspace{0pt}
\centering
\begin{tikzpicture}[ultra thick]
 \draw [black,   -latex      ] (0,0.5) -- (0,0) node [] {};
\end{tikzpicture}
\end{minipage}

\vspace{0pt}
\begin{subfigure}[b]{\linewidth}
\centering
\begin{minipage}{.85\textwidth}
\begin{lstlisting}[mathescape,numbers=left]
int4 mul3(int4 x) { return x * 3; }
kernel vectorScal(global int* in,out, int len){
 for (int i=get_group_id; i < len/1024;
      i+=get_num_groups) {
  global int* grp_in  = in+(i*1024);
  global int* grp_out = out+(i*1024);
  for (int j=get_local_id; j < 1024/4;
       j+=get_local_size) {
    global int4* in_vec4 =(int4*)grp_in+(j*4);
    global int4* out_vec4=(int4*)grp_out+(j*4);
    *out_vec4 = mul3(*in_vec4);      
} } }  
\end{lstlisting}
\end{minipage}
\caption{\textbf{OpenCL program} produced by our code generator.}
\label{fig:codeex:ocl}
\end{subfigure}
\vspace{-25pt}
\caption{
Pseudo-code representing vector scaling.
The user maps the \code{mul3} function over the input array~(\subref{fig:codeex:map}).
This high-level expression is transformed into a low-level expression~(\subref{fig:codeex:impl}) using rewrite rules.
Finally, our code generator turns the low-level expression into an OpenCL program~(\subref{fig:codeex:ocl}).
}
  \label{fig:codeex}
\end{figure}

We illustrate the advantages of our approach using a simple vector scaling example shown in Figure~\ref{fig:codeex}.
The user expresses the computation by writing a high-level expression using our \pat{map} pattern as shown in Figure~\ref{fig:codeex:map}.
This coding style is similar to functional and dataflow programming.

Our technique first rewrites the high-level expression into something closer to the OpenCL programming model.
This is achieved by applying the rewrite rules presented later in Section~\ref{rules}.
Figure~\ref{fig:codeex:impl} shows one possible derivation of the original high-level expression.
The $\circ$ operator represents sequential function composition. %, \ie $f \circ g(x) = f(g(x))$.
Starting from the last line, the input is split into chunks of 1024 elements.
Each chunk is mapped onto a group of threads, called \emph{workgroup}, with the \pat{map-workgroup} low-level pattern (line~2).
Within a workgroup (lines~3--5), we vectorize the elements (line~5), each mapped to a local thread inside a workgroup via the \pat{map-local} pattern (line~3).
Each local thread now processes 4 elements, enclosed in a vector type.
Finally, the \pat{vect-4} pattern (line~4) vectorizes the user defined function \code{mul3}.
The exact meaning of our patterns will be given later in Section~\ref{pattern}.

The last step consists of traversing the low-level expression and generating OpenCL code for each low-level pattern encountered (Figure~\ref{fig:codeex:ocl}).
The two map patterns generate the for-loops (line~3--4 and~7--8) that iterate over the input array assigning work to the workgroups and local threads.
The information of how many chunks each workgroup and thread processes comes from the corresponding \pat{split}.
In line~11 the vectorized version of the user defined \code{mul3} function (defined in line~1) is finally applied to the input array.

To summarize, our approach is able to generate OpenCL code starting from a high-level program representation. % of a program.
This is achieved by systematically lowering the high-level expression into a low-level form suitable for code generation.
The next two sections present our high-level and low-level patterns, the code generation mechanism and the rewrite rules. % in more details.
\from{PLDI end}



\section{Patterns: Design and Implementation}
\label{section:patterns}

In this section, we will introduce the \emph{patterns} which form the expressions used for code generation.
As we saw in the previous section there exist two type of patterns: high-level algorithmic patterns and low-level \OpenCL patterns.
Some of the high-level algorithmic patterns directly correspond to algorithmic skeletons we have introduced in \autoref{chapter:skelcl}.
As we will also introduce patterns which we have not seen so far and which do not oblige to the common definition of algorithmic skeletons, we will use the more generic term \emph{pattern} throughout this and the next chapter.

The key idea of our approach is to expose algorithmic choices and hardware-specific program optimizations as patterns that are systematically derived using a rule rewriting system (discussed later in \autoref{section:rules}).
The high-level algorithmic patterns represent structured parallelism.
They can either be used by the programmer directly as a stand-alone language, as a domain specific language embedded in a general purpose programming language, or used as an intermediate representation targeted by the compiler of another language.
Once a program is represented with our high-level patterns, we automatically transform the program into low-level patterns.
The low-level patterns represent hardware-specific concepts expressed by a programming model such as \OpenCL, which is the target chosen for this thesis.
Following the same approach, a different set of low-level patterns could be designed to target other low-level programming models such as Pthreads or MPI.


\subsection{High-level Algorithmic Patterns}

We define our patterns as functions.
To simplify our implementation, we encode all types as arrays with primitives represented by arrays of length 1.
The only exceptions are the user-defined functions, such as the \code{mul3} function in \autoref{fig:codeex:map} that operates on a primitive type.

\autoref{tab:hlskel} presents our high-level patterns used to define programs at the algorithmic level.
Most of the patterns are well known in functional programming, like \map and \reduce.
The \zip, \splitN and \join patterns transform the shape of the data.
The \iterateN pattern iteratively applies a function multiple times.
Finally, the \reorder pattern lets our system know that it is safe to reorder the elements of an array arbitrarily, which enables additional optimizations -- as we will see later in \autoref{chapter:codeGeneration-evaluation}.

In the following, we discuss each high-level algorithmic pattern in more detail including their formal definitions.
As in \autoref{chapter:skelcl} we use the Bird-Meertens formalism~\cite{Bird88} as inspiration for our notation of the patterns.
For full details see~\autoref{par:notation} on page~\pageref{par:notation}.
Here is a short reminder of our notation:
we write function application using a space and functions are often curried;
binary operators (\eg, $\oplus$) are written using infix notation and can be written using prefix notation when being sectioned by using parenthesis: $(\oplus)$.
For an array $xs$ of length $n$ with elements $x_i$ we write $[x_1, \ldots, x_n]$.

In this chapter we are especially interested in how patterns can be composed and nested.
As \emph{types} formally specify which compositions and nesting of patterns are legal, we will define the type of each pattern.
% For expressing types we base our syntax on the syntax established by J. Roger Hindley, Robin Miler, and Luis Damas as part of the Hindley---Milner type system~\cite{}.
We write $e : \sigma$ to denote that expression $e$ has type $\sigma$.
% To make typing judgments we write $e_0 : \sigma_0,\ \ldots,\ e_n : \sigma_n\ \vdash e : \sigma$ to denote that under the assumptions that $e_0,\ \ldots,\ e_n$ have types $\sigma_0,\ \ldots,\ \sigma_n$ the expression $e$ has type $\sigma$.
For a function mapping values of type $\alpha$ to values of type $\beta$ we write its type as $(\alpha \rightarrow \beta)$.
Tuple types are written as $\langle\alpha, \beta\rangle$.
Finally, arrays have their length denoted as part of their type:
for an array with elements of type $\alpha$ and length $n$, we write $[\alpha]_n$.

A more formal description of a core subset of the patterns described in this section can be found in~\cite{SteuwerFeLiDu2015}, where formal syntax, typing rules, and formal semantics are given.

\begin{table}[t]
\centering
\begin{tabular}{p{.1\textwidth}p{.85\textwidth}}
\toprule
\tabhead{Pattern} & \tabhead{Description}\\
\midrule
 \map
     & Apply a given function to every element of an input array.\\
 \reduce
     & Perform a reduction of an input array using a user-defined binary function and an initial value.\\
 \zip
     & Build an array of pairs by pairwise combining two arrays.\\
 \splitN
     & Produce a multi-dimensional array by splitting an array in chunks of a given size.\\
 \join
     & Join the two outer most dimensions of an multi-dimensional array.\\
 \iterateN
     & Iterate a given function over an input array a fixed number of times.\\
 \reorder
     & Reorder the element of the input array.\\
\bottomrule
\end{tabular}
\caption{High-level algorithmic patterns used by the programmer.}
\label{tab:hlskel}
\end{table}


\paragraph{Map}
The \map pattern is well known in functional programming and applies a given unary function $f$ to all elements of an input array.
In \autoref{chapter:skelcl}, we defined the \map pattern as an algorithmic skeleton (see \autoref{definition:map}).
The same definition holds here.
We repeat it here for completeness and we add the type information:
\begin{definition}
  \label{definition:pattern:map}
  Let $xs$ be an array of size $n$ with elements $x_i$ where $0 < i \leq n$.
  Let $f$ be a unary customizing function defined on elements.
  The \map pattern is then defined as follows:
  \begin{equation*}
    \map\ f\ [x_1, x_2, \dots, x_n] \eqdef [f\ x_1, f\ x_2, \dots, f\ x_n]
  \end{equation*}
  The types of $f$, $xs$, and \map are as follows:
  \begin{align*}
    f &: (\alpha \rightarrow \beta),\\
    xs &: [\alpha]_n,\\
    \map\ f\ xs &: [\beta]_n
  \end{align*}
\end{definition}

\noindent
In \autoref{chapter:skelcl} we also defined the \map skeleton for operating on matrices (see \autoref{definition:map:matrix}).
In this chapter we represent matrices as nested arrays, therefore, performing an operation on each element of a matrix can be represented by nesting two \map patterns:
\begin{align*}
  mapMatrix\ \ f\ \ X &= \map\ (\map\ f)\ X
\end{align*}
Let us assume that $X$ represents an $n\times m$ matrix with elements of type $\alpha$, then its type is $\big[[\alpha]_m\big]_n$.
The outer \map applies its customizing function to every row of matrix $X$.
The customizing function is defined by currying \map and $f$, thus, producing a function which applies $f$ to every element of its argument array.
Therefore, $f$ will be applied to every element of matrix $X$.

\paragraph{Reduce}
The \reduce pattern (\aka, fold or accumulate) uses a binary operator $\oplus$ to combine all elements of the input array.
We require the operator $\oplus$ to be associative and commutative which allows for an efficient parallel implementation.
By requiring commutativity, our system can also generate vectorized implementations of the reduction and utilize the efficient coalesced memory access pattern, as we saw in \autoref{section:reduce:case-study}.
In \autoref{chapter:skelcl} we defined \reduce as an algorithmic skeleton (see \autoref{definition:reduce}).
The same definition holds here.
We repeat it here for completeness and we add the type information:
\begin{definition}
  \label{definition:pattern:reduce}
  Let $xs$ be an array of size $n$ with elements $x_i$ where $0 < i \leq n$.
  Let $\oplus$ be an associative and commutative binary customizing operator with the identity element $\id_\oplus$.
  The \reduce pattern is then defined as follows:
  \begin{equation*}
    \reduce\ (\oplus)\ \id_\oplus\ [x_1, x_2, \dots, x_n]
      \eqdef [x_1 \oplus x_2 \oplus \dots \oplus x_n]
  \end{equation*}
  The types of $(\oplus)$, $\id_\oplus$, $xs$, and \reduce are as follows:
  \begin{align*}
    (\oplus) &: ((\alpha, \alpha) \rightarrow \alpha),\\
    \id_\oplus &: \alpha,\\
    xs &: [\alpha]_n,\\
    \reduce\ (\oplus)\ \id_\oplus\ xs &: [\alpha]_1
  \end{align*}
\end{definition}


\paragraph{Zip}
The \zip pattern and the following \splitN/\join patterns transform the shape of the data. %which we store in the type system, \ie, number of dimensions and size of each dimension.
The \zip pattern fuses two arrays into a single array of pairs.

\begin{definition}
  \label{definition:pattern:zip}
  Let $xs$ and $ys$ be arrays of size $n$ with elements $x_i$ and $y_i$ where $0 < i \leq n$.
  The \zip pattern is then defined as follows:
  \begin{equation*}
    \zip\ [x_1, x_2, \dots, x_n]\ [y_1, y_2, \dots, y_n]\\
      \eqdef [\langle x_1, y_1\rangle, \langle x_2, y_2\rangle, \dots, \langle x_n, y_n\rangle]
  \end{equation*}
  The types of $xs$, $ys$, and \zip are as follows:
  \begin{align*}
    xs &: [\alpha]_n,\\
    ys &: [\beta]_n,\\
    \zip\ xs\ ys &: [\langle\alpha, \beta\rangle]_n
  \end{align*}
\end{definition}

\noindent
This definition significantly differs from the definition of the \zip skeleton in \autoref{chapter:skelcl}:
While in \autoref{definition:zip}, \zip applies a given function to a pair of elements, there is no function to be applied in \autoref{definition:pattern:zip} of the \zip pattern.
The behavior of the \zip skeleton can be achieved by composing the \zip pattern with the \map pattern:
\begin{align*}
  zipWith\ f\ xs\ ys &= \map\ f\ (\zip\ xs\ ys)
\end{align*}


\paragraph{Split and Join}
The \splitN pattern, which is most often combined with a \join, partitions an array into chunks of specific size, resulting in an extra dimension in the output array.

We start with the definition of the \splitN pattern.
\begin{definition}
  \label{definition:pattern:split}
  Let $n$ be an integer value.
  Let $xs$ be an array of size $m$ with elements $x_i$ where $0 < i \leq m$.
  Let us assume that $m$ is evenly divisible by $n$.
  The \splitN pattern is then defined as follows:
  \begin{align*}
    &\splitN\ n\ [x_1, x_2, \dots, x_m] \eqdef\\
    &\qquad\big[[x_1, \ldots, x_n], [x_{n+1}, \ldots, x_{2n}], \ldots, [x_{m-n+1}, \ldots, x_m]\big]
  \end{align*}
  The types of $n$, $xs$, and \splitN are as follows:
  \begin{align*}
    n &: int,\\
    xs &: [\alpha]_m,\\
    \splitN\ n\ xs &: \big[[\alpha]_n\big]_{\frac{m}{n}}
  \end{align*}
\end{definition}

\bigskip
%The formal type of the \splitN pattern is:
%\begin{align}
%  n : int,\ xs : [\alpha]_m\ \vdash\ split\ n\ xs : \big[[\alpha]_n\big]_{{}^m/_n}
%\end{align}

\noindent
The corresponding \join pattern does the opposite:
it reassembles an array of arrays by merging dimensions.
\begin{definition}
  \label{definition:pattern:join}
  Let $xs$ be an array of size $\frac{m}{n}$ whose elements are arrays of size $n$.
  We denote the elements of the $i$th inner array as $x_{((i-1)\times n) + j}$ where $0 < i \leq \frac{m}{n}$ and $0 < j \leq n$.
  The \join pattern is then defined as follows:
  \begin{align*}
    &\join_n\ \big[[x_1, \ldots, x_n], [x_{n+1}, \ldots, x_{2n}], \ldots, [x_{m-n+1}, \ldots, x_m]\big] \eqdef\\
    &\qquad[x_1, x_2, \dots, x_m]
  \end{align*}
  The types of $xs$ and $\join_n$ are as follows:
  \begin{align*}
    xs &: \big[[\alpha]_n\big]_{\frac{m}{n}},\\
    \join_n\ xs &: [\alpha]_m
  \end{align*}
\end{definition}

\noindent
We will almost always omit the subscript, as $n$ can be inferred from the input array.
From these definitions it follows, that the compositions of $\splitN\ n$ and $\join_n$: $\join_n \circ \splitN\ n$ and $\splitN\ n \circ \join_n$ for any value $n$ yields the same type and also does not change any element in the array, \ie, it is equivalent to the identify function \emph{id}.

These two patterns used together are similar to the split-join concept from data-flow languages such as StreamIt~\cite{ThiesKaAm2002}.


\paragraph{Iterate}
The \iterateN pattern corresponds to the mathematical definition of iteratively applying a function, which is defined as: {$f^0 = id$} and {$f^{n+1} = f^n \circ f$}.

\begin{definition}
  \label{definition:pattern:iterate}
  Let $n$ be an integer value with $n \geq 0$.
  Let $f$ be a unary function on arrays.
  Let $xs$ be an array of arbitrary size.
  We define the \iterateN pattern recursively:
  \begin{align*}
    \iterateN\ 0\ f\ xs &\eqdef xs,\\
    \iterateN\ n\ f\ xs &\eqdef \iterateN\ (n-1)\ f\ (f\ xs)
  \end{align*}
  The types of $n$, $f$, $xs$, and \iterateN are as follows:
  \begin{align*}
    n &: int,\\
    f &: ([\alpha]_k \rightarrow [\alpha]_{F(k)}),
      \begin{aligned}[t]
        & \forall k \text{ and where } \\
        &F : (int\rightarrow int) \text{ describes the change}\\
        &\text{of array length when applying } f,
      \end{aligned}\\
    xs &: [\alpha]_m,\\
    \iterateN\ n\ f\ xs &: [\alpha]_{F^n(m)}
  \end{align*}
\end{definition}


% In terms of implementation, our code generator emits a for-loop to perform the iteration, and two pointers for input and output.
% After each iteration, we swap the pointers, so that the output of the last iteration becomes the input for the next one.
% We automatically allocated and calculate memory requirements based on the information from the input and output type.
\noindent
The type of the \iterateN pattern is interesting as its result type depends on the effect $f$ has on the size of its argument.
The index function $F$ describes the effect the function $f$ has on the length of its input array.
This index function is used to compute the effect the \iterateN function has when applying $f$ $n$ times on its input array.
Please note, that the length of the input array of $f$, \ie, $k$, possibly changes every time $f$ is applied by \iterateN.
A more formal treatment of the topic can be found in~\cite{SteuwerFeLiDu2015}.

%\begin{align}
%  n : int,\ F : (int \rightarrow int),\ f : (\forall k\ .\ [\alpha]_k \rightarrow [\alpha]_{F(k)}),\ xs : [\alpha]_m\ %
%  \vdash iterate\ n\ f\ xs : [\alpha]_{F^{n}(m)}
%\end{align}

\paragraph{Reorder}
The \reorder pattern is used to specify that the ordering of the elements of an array does not matter.

\begin{definition}
  \label{definition:pattern:reorder}
  Let $xs$ be an array of size $n$ with elements $x_i$ where $0 < i \leq n$.
  Let $\sigma$ be an arbitrary permutation of $[1,\ldots, n]$.
  The \reorder pattern is then defined as follows:
  \begin{align*}
    \reorder\ [x_1, \ldots, x_n] &\eqdef [x_{\sigma(1)}, \ldots, x_{\sigma(n)}]
  \end{align*}
  The types of $xs$ and \reorder are as follows:
  \begin{align*}
    xs &: [\alpha]_n,\\
    \reorder\ xs &: [\alpha]_n
  \end{align*}
\end{definition}

\noindent
This definition allows to pick any permutation to reorder the elements of an array arbitrarily which may enable optimizations, as we will see later.





\subsection{Low-level, \OpenCL-specific Patterns}

% Programming manycore CPU and GPU devices is difficult due to the need for managing parallelism, the memory hierarchy and other hardware specific features.
In order to achieve the highest performance, programmers often use a set of intuitive ``rules of thumb'' to optimize their applications.
We extensively discussed one application example in \autoref{section:reduce:case-study}.
Each hardware vendor provides own optimization guides~\cite{CUDAProgrammingGuide,AMDProgrammingGuide,IntelGPUProgrammingGuide,IntelXeonProgrammingGuide} that extensively cover vendor's hardware particularities and optimizations.
The main idea behind our approach is to identify common optimization patterns and express them systematically rather than intuitively, using low-level patterns coupled with a rewrite-rule system.

\autoref{tab:llskel} gives an overview of the \OpenCL-specific patterns we have identified.

\begin{table}[t]
\centering
\begin{tabular}{p{.25\textwidth}p{.7\textwidth}}
\toprule
\tabhead{Pattern} & \tabhead{Description}\\
\midrule
 \mapWorkgroup
     & Each \OpenCL \textbf{work-group} applies the given function to an element of the input array.\\
 \mapLocal
     & Each \textbf{local work-item} of a work-group applies the given function to an element of the input array.\\
 \mapGlobal
     & Each \textbf{global work-item} applies the given function to an element of the input array.\\
 \mapWarp
     & Each \textbf{warp} applies the given function to an element of the input array.\newline Only available for Nvidia \GPUs.\\
 \mapLane
     & Each \textbf{work item inside a warp} applies the given function to an element of the input array.\newline Only available for Nvidia \GPUs.\\
 \mapSeq
      & Apply the given function to every element of the input array \textbf{sequentially}.\\
 \reduceSeq
      & Perform the reduction using the given binary reduction function and initial value on the input array \textbf{sequentially}.\\
 \reorderStride
      & Access input array with a given stride to maintain \textbf{memory coalescing}.\\
 \toLocal
      & Store the results of a given function to \textbf{local memory}.\\
 \toGlobal
      & Store the results of a given function to \textbf{global memory}.\\
 \asVector
      & Turns the elements of an array into \textbf{vector type} of a given width.\\
 \asScalar
      & Turns the elements of an array into \textbf{scalar type}.\\
 \vect
      & \textbf{Vectorize} a given function by a given width.\\
\bottomrule
\end{tabular}
\caption{Low-level \OpenCL patterns used for code generation.}
\label{tab:llskel}
\end{table}

\paragraph{Parallel Maps}
The different low-level \OpenCL \map patterns represent possible ways of mapping computations to the hardware and exploit thread-level parallelism in \OpenCL.
The execution semantics and types of all these low-level \OpenCL \map patterns are the same as of the high-level \map pattern shown in \autoref{definition:pattern:map}.
In fact, these patterns can be viewed as \emph{instantiations} of the \emph{interface} the high-level \map pattern defines, as the low-level \OpenCL pattern provide different possibilities to implement the \map pattern.

\todo{sg: hier waere ein Bild gut, das diese maps zusammen mit der OpenCL Thread-Hierarchie zeigt}
The \mapWorkgroup pattern assigns work to an \OpenCL work-group, \ie, following \autoref{definition:pattern:map} each \OpenCL work-group executes the given function on a different part of the input vector.

The \mapLocal pattern assigns work to a local work-item inside a work-group.
This pattern only makes sense in the context of a work-group and is, therefore, only correctly used when nested inside a \mapWorkgroup pattern, \eg, $\mapWorkgroup\ (\mapLocal\ f)$.

The \mapGlobal pattern can be used to assign work to work-items outside of a work-group.
This allows us to map computations in different ways to the thread hierarchy of \OpenCL.

There are two additional patterns which are only valid when generating code for Nvidia \GPUs:
\mapWarp and \mapLane.
These two patterns capture the notion of \emph{warps} present in Nvidia's \GPU architectures, \ie, 32 work-items are grouped and executed together (see \autoref{chapter:background} for details).
Barrier synchronizations between work-items in the same warp can be avoided because a warp executes in a lock step manner.
To exploit this optimization we provide the \mapWarp pattern which assigns work to a warp, \ie, a group of 32 work-items.
The \mapLane pattern is used to assign work to an individual work-item inside a warp.

%The code generation for all these map patterns is similar; we describe it using \pat{map-workgroup} as an example.
%A loop is generated, where the iteration variable is determined by the \emph{workgroup-id} function provided by \OpenCL.
%Inside of the loop, a pointer is generated to partition the input array, so that every work-group calls \pat{map-workgroup}'s customizing function on a different chunk of data.
%An output pointer is generated similarly.
%We continue with the body of the loop by generating the code for the customizing function recursively.
%Finally, an appropriate synchronization mechanism for the given map pattern is added.
%For instance after a \pat{map-local} we add a barrier synchronization to synchronize the threads inside of the work-group.

\FloatBarrier

\paragraph{Sequential Map and Reduce}
The \mapSeq and \reduceSeq patterns perform a sequential map and reduction, respectively, within a single work-item.
% In both cases the generated code consists of a loop iterating over the array and calling the customizing function.
% For the reduction an accumulation variable is initialized with the given initial value and used to accumulate the results produced by the successive calls to to the customizing function.

For the \mapSeq pattern, the semantics and type are the same as for the high-level \map pattern shown in \autoref{definition:pattern:map}.
This is not the case for the \reduceSeq and \reduce patterns, where their types differ.
For the high-level \reduce pattern, we require the customizing binary operator to be associative and commutative in order to allow for an efficient parallel implementation.
As the \reduceSeq pattern performs a sequential reduction, we can relax these requirements, therefore, we define \reduceSeq separately.
\begin{definition}
  \label{definition:pattern:reduceSeq}
  Let $xs$ be an array of size $n$ with elements $x_i$ where $0 < i \leq n$.
  Let $\oplus$ be a binary customizing operator with the identity element $\id_\oplus$.
  The \reduceSeq pattern is then defined as follows:
  \begin{equation*}
    \reduceSeq\ (\oplus)\ \id_\oplus\ [x_1, x_2, \dots, x_n]
      \eqdef [(\dots((id_\oplus \oplus x_1) \oplus x_2) \ldots \oplus x_n)]
  \end{equation*}
  The types of $(\oplus)$, $\id_\oplus$, $\vec{x}$, and \reduce are as follows:
  \begin{align*}
    (\oplus) &: ((\alpha, \beta) \rightarrow \alpha),\\
    \id_\oplus &: \alpha,\\
    xs &: [\beta]_n,\\
    \reduceSeq\ (\oplus)\ \id_\oplus\ xs &: [\alpha]_1
  \end{align*}
\end{definition}


\paragraph{Reorder-stride}
The \reorderStride pattern enforces a special reordering of an array.
The produced strided-access pattern ensures that after splitting the workload, consecutive work-items access consecutive memory elements.
This corresponds to the \emph{coalesce memory accesses}, which are beneficial on modern \GPUs as discussed in \autoref{chapter:background}.
\begin{definition}
  \label{definition:pattern:reorderStride}
  Let $s$ be an integer value.
  Let $xs$ be an array of size $m$ with elements $x_i$, where $0 < i \leq m$.
  Let us assume that $m$ is evenly divisible by $s$ and that $m = s\times n$ for some integer value $n$.
  The \reorderStride pattern is then defined as follows:
  \begin{align*}
    \reorderStride\ s\ [x_1, x_2, \dots, x_m] &\eqdef [y_1, y_2, \dots, y_m], \text{ where}\\
    y_i &\eqdef x_{i\ \text{\normalfont div } n + s \times (i \bmod{n})}
  \end{align*}
  Where $x\ \text{\normalfont div}\ y$ is integer division and $x \bmod{y}$ is the modulo operation.
  The types of $s$, $xs$, and \reorderStride are as follows:
  \begin{align*}
    s &: int,\\
    xs &: [\alpha]_m,\\
    \reorderStride\ s\ xs &: [\alpha]_m
  \end{align*}
\end{definition}

Our code generator produces no code directly, but rather the reordering determines how the data is read from the array in the following pattern.

%\noindent
\todo{Maybe a figure visualizing the reordering ...}

%Our implementation does not produce code directly, but generates instead an index function, which is used when accessing the array the next time.
%Therefore, any following read implicitly reorders the array.
%Our design is general enough to supports user-defined index functions as well, which we will add in the future.
%The type of \pat{reorder-stride} corresponds to the type of the high-level \pat{reorder} pattern.

\paragraph{{\footnotesize to}Local and {\footnotesize to}Global}
The \toLocal and \toGlobal patterns are used to specify where the result of a given function $f$ should be stored.
As explained in more detail in \autoref{chapter:background}, \OpenCL defines two distinct address spaces: global and local.
Global memory is the commonly used large but slow memory.
On \GPUs, the comparatively small local memory has a high bandwidth with low latency and is used to store frequently accessed data.
With these two patterns, we can in effect exploit the memory hierarchy defined in \OpenCL.

First, we define \toLocal:
\begin{definition}
  \label{definition:pattern:toLocal}
  Let $f$ be a function.
  The \toLocal pattern is then defined as follows:
  \begin{align*}
    \toLocal\ f &\eqdef f',
    \begin{aligned}[t]
      &\text{ where $f'\ x \eqdef f\ x$, $\forall x$, and  $f'$ is guaranteed to store}\\
      &\text{its result in local memory.}
    \end{aligned}
  \end{align*}
  The types of $f$, and \toLocal are as follows:
  \begin{align*}
    f &: (\alpha \rightarrow \beta)\\
    \toLocal\ f &: (\alpha \rightarrow \beta)
  \end{align*}
\end{definition}

\noindent
The definition of \toGlobal is correspondent:
\begin{definition}
  \label{definition:pattern:toGlobal}
  Let $f$ be a function.
  The \toGlobal pattern is then defined as follows:
  \begin{align*}
    \toGlobal\ f &\eqdef f',
    \begin{aligned}[t]
      &\text{ where $f'\ x \eqdef f\ x$, $\forall x$ and $f'$ is guaranteed to store}\\
      &\text{ its result in global memory.}
    \end{aligned}
  \end{align*}
  The types of $f$, and \toGlobal are as follows:
  \begin{align*}
    f &: (\alpha \rightarrow \beta)\\
    \toGlobal\ f &: (\alpha \rightarrow \beta)
  \end{align*}
\end{definition}

% These patterns act similarly to a typecast and are in fact implemented as such so that no code is emitted directly.

%In our design, every function reads its input and writes its output using pointers provided by the callee function.
%As a result, we can simply force a store to local memory by wrapping any function with our \pat{toLocal} pattern.
%In the code generator, this will simply change the output pointer of function $f$ to an area in local memory.

%The types of \pat{toLocal} and \pat{toGlobal} are identical and as follows:
%\begin{align}
%  f : (\alpha \rightarrow \beta)\ &\vdash\ toLocal\ f : (\alpha \rightarrow \beta)
%  \label{eq:type:toLocal}
%  \\
%  f : (\alpha \rightarrow \beta)\ &\vdash\ toGlobal\ f : (\alpha \rightarrow \beta)
%  \label{eq:type:toGlobal}
%\end{align}


\paragraph{{\footnotesize as}Vector, {\footnotesize as}Scalar, and Vectorize}
The \OpenCL programming model supports vector data types such as \code{float4} where any operations on this type will be executed in the hardware vector units.
In the absence of vector units in the hardware, the \OpenCL compiler generates a version of the code solely using scalar data types automatically.

The \asVector and \asScalar patterns change the data type into vector elements and scalar elements, correspondingly.
For instance, an array of \code{float} is transformed into an array of \code{float4} as seen in the motivation example (\autoref{fig:codeex}).
The \vect pattern vectorizes the given function by simply converting all the operations that apply to vector types into vectorized operations.
%Our current implementation can only vectorize functions containing simple arithmetic operations such as $+$ or $-$.
In case of more complex functions, we rely on external tools~\cite{KarrenbergHa2011} for vectorizing the operations.
%These tools are not required to performing further analysis to find opportunities for vectorization.
%The rewrite rules presented in \autoref{section:rules} ensure that vectorization is only applied to patterns where vectorization is a legal optimization.

We start by defining the \asVector pattern.
\begin{definition}
  \label{definition:pattern:asVector}
  Let $n$ be a positive integer value.
  Let $xs$ be an array of size $m$ with elements $x_i$ where $0 < i \leq m$.
  Let us assume, that $m$ is evenly divisible by $n$.
  The \asVector pattern is then defined as follows:
  \begin{align*}
    &\asVector\ n\ [x_1, x_2, \dots, x_m] \eqdef \\
    &\qquad\big[\{x_1, \ldots, x_n\}, \{x_{n+1}, \ldots, x_{2n}\}, \ldots, \{x_{m-n+1}, \ldots, x_m\}\big],
  \end{align*}
  where $\{x_1,\ldots,x_n\}$ denotes a vector of width $n$.\\
  The types of $n$, $xs$, and \asVector are as follows:
  \begin{align*}
    n &: int\\
    xs &: [\alpha]_m,\\
    \asVector\ n\ xs &: [\alpha_n]_{\frac{m}{n}}
  \end{align*}
  Here $\alpha$ is required to be a basic scalar type, \eg, \code{int} or \code{float}, and $\alpha_n$ denotes the vectorized version of that type with a vector width of $n$.
\end{definition}

\noindent
The corresponding \asScalar pattern is defined as follows.
\begin{definition}
  \label{definition:pattern:asScalar}
  Let $xs$ be an array of size $\frac{m}{n}$ whose elements are vectors of width $n$.
  We denote the individual vector elements of the $i$th element of $xs$ as $x_{((i-1)\times n)+j}$ where $0 < i \leq \frac{m}{n}$ and $0 < j \leq n$.
  The \asScalar pattern is then defined as follows:
  \begin{align*}
    &\asScalar_n\ \big[\{x_1, \ldots, x_n\}, \{x_{n+1}, \ldots, x_{2n}\}, \ldots, \{x_{m-n+1}, \ldots, x_m\}\big] \eqdef\\
    &\qquad [x_1, x_2, \dots, x_m],
  \end{align*}
  where $\{x_1,\ldots,x_n\}$ denotes a vector of width $n$.\\
  The types of $xs$, and $\asScalar_n$ are as follows:
  \begin{align*}
    xs &: [\alpha_n]_{\frac{m}{n}},\\
    \asScalar_n\ xs &: [\alpha]_m
  \end{align*}
  Here $\alpha$ is required to be a basic scalar type, \eg, \code{int} or \code{float}, and $\alpha_n$ denotes the vectorized version of that type with a vector width of $n$.
\end{definition}

We will almost always omit the subscript, as $n$ can be inferred from the input array.

\noindent
Finally, we define the \vect pattern.
\begin{definition}
  \label{definition:pattern:vect}
  Let $n$ be a positive integer value.
  Let $f$ be a function.
  The \vect pattern is then defined as follows:
  \begin{align*}
    \vect\ n\ f &\eqdef f_n, \text{ where } f_n\ \{x_1, \ldots, x_n\} = \{f\ x_1, \ldots, f\ x_n \}
  \end{align*}
  and $\{x_1,\ldots,x_n\}$ denotes a vector of width $n$.\\
  The types of $n$, $f$, and \vect are as follows:
  \begin{align*}
    n &: int,\\
    f &: (\alpha \rightarrow \beta),\\
    \vect\ n\ f &: (\alpha_n \rightarrow \beta_n)
  \end{align*}
  Here $\alpha$ and $\beta$ are required to be a basic scalar type, \eg, \code{int} or \code{float}, and $\alpha_n$ denotes the vectorized version of that type with a vector width of $n$.
\end{definition}

\subsection{Summary}
In this section we introduced two type of \emph{patterns}:
high-level algorithmic patterns and low-level, \OpenCL-specific patterns.
While all of these patterns can be used by the application programmer to describe the solution for a particular problem, we expect the programmer to focus on the algorithmic patterns which should be used to express a high-level algorithmic implementation of the problem solution.
We will see in the next section how such an implementation composed of our high-level algorithmic patterns can be systematically rewritten using \emph{rewrite rules}.
During this process the original implementation will be modified and \OpenCL-specific patterns will be used.




\section{Rewrite Rules}
\label{section:rules}

This section introduces our set of rewrite rules that transform high-level expressions written using our algorithmic patterns into semantically equivalent expressions.
One goal of our approach is to keep each rule as simple as possible and only express one fundamental concept at a time.
For instance the vectorization rule, as we will see, is the only place where we express the vectorization concept.
This is different from most prior approaches that would produce a special vectorized version of different algorithmic patterns such as map or reduce.
The superiority of our approach lies in the power of composition;
many rules can be applied successively to produce expressions that compose hardware concepts or optimizations and that are provably correct by construction.

Similarly to our patterns, we distinguish between algorithmic and lowering rules.
Algorithmic rules produce derivations that represent the different algorithmic choices.
Our \OpenCL-specific rules map expressions to \OpenCL patterns.
Once the expression is in its lowest form, it is possible to produce \OpenCL code for each single pattern easily with our code generator as described in the previous section.


%\paragraph{Syntax and Rule Derivation}
%Some rules can only be activated if certain conditions are true.
%We use the syntax $[pre:post]$ to represent the pre and post conditions of a rule.
%The pre condition $pre$ corresponds to the list of conditions that must be true for the rule to be applied.
%The post condition $post$ is set for any function bound to the pattern.
%The $\overline{\rule[-.3\baselineskip]{0pt}{1.5ex}\hspace{1.5ex}}$, $\wedge$, and $\vee$ corresponds to the logical \emph{not}, \emph{and}, and \emph{or} operators respectively.

%We use leftmost derivation when applying the rules, which means that the leftmost non-terminal is always derived first.

\newenvironment{rerule}[1]%
{\begin{equation}\begin{array}{#1}\ignorespaces}%
{\end{array}\end{equation}%
\ignorespacesafterend}

\newenvironment{rerule*}[1]%
{\begin{equation*}\begin{array}{#1}\ignorespaces}%
{\end{array}\end{equation*}%
\ignorespacesafterend}



\subsection{Algorithmic Rules}

\paragraph{Identity}
The identity rule in \autoref{eq:algo:identity} specifies that it is always valid to compose any function with the identity function \pat{id}.
The \pat{id} function can acts as a copy operation; this is useful to express copy to local memory for instance when composed with the corresponding lowering rule.
%
\begin{rerule}{lclcl}
  f & \rightarrow & f \circ \pat{id} & | & \pat{id} \circ f
  \label{eq:algo:identity}
\end{rerule}

 
\paragraph{Iterate decomposition}
The rule in \autoref{eq:algo:iterate} expresses the fact that an iteration can be decomposed into several iterations.
%
\begin{rerule}{lcl}
  \pat{iterate}^{m+n}\ f
    & \rightarrow &
      \pat{iterate}^m\ f
        \circ \pat{iterate}^n\ f
  \label{eq:algo:iterate}
\end{rerule}

\paragraph{Reorder commutativity}
\autoref{eq:algo:reorder} shows that if the data can be reordered arbitrarily it does not matter if we apply a function $f$ to each element before or after the reordering.
%
\begin{rerule}{lcl}
  \pat{map}\ f \circ \pat{reorder}
    & \rightarrow & \pat{reorder} \circ \pat{map}\ f\\
  \pat{reorder} \circ \pat{map}\ f
    & \rightarrow & \pat{map}\ f \circ \pat{reorder}  
  \label{eq:algo:reorder}
\end{rerule}

\paragraph{Split-join}
The split-join rule in \autoref{eq:algo:splitjoin} partitions a map into two maps.
This allows us to nest map patterns in each other and, thus, \emph{maps} the computation to the thread hierarchy of the \OpenCL programming model, such as \pat{map-workgroup\ (map-local\ $f$)} as seen in our motivation example (\autoref{fig:codeex}).
%
\begin{rerule}{lcl}
  \pat{map}\ f
    & \rightarrow &
      \pat{join}
        \circ \pat{map\ (map\ $f$)}
        \circ \pat{split}\ n
  \label{eq:algo:splitjoin}
\end{rerule}


\paragraph{Reduction}
The reduction in \autoref{eq:algo:red} is currently our most complex rule but also the most powerful one.
It expresses the reduction function as a composition of other primitive functions, which is a fundamental aspect of our work.
From the algorithmic point of view we first define a partial reduction pattern \pat{part-red}.
This partial reduction reduces an array of $n$ elements to an array of $m$ elements where $1 \leq m < n$.
The reduction can be derived in a partial reduction combined with a full reduction which ensures we end up with one unique element.
% Another possible derivation consists in iterating a partial reduction until a full reduction is achieved (this is what $\infty$ represents).
%Note that our definition of \pat{reduce} remains correct since the result of a partial reduction is always composed with the reduction to ensure we end up with one unique element.
%
\begin{rerule}{lcl}
  \pat{reduce}\ f\ z
    & \rightarrow &
      \pat{reduce}\ f\ z \circ \pat{part-red}\ f\ z
  \label{eq:algo:red}
\end{rerule}

\paragraph{Partial Reduction}
The first possible derivation for partial reduction, in \autoref{eq:algo:part-red}, leads to the full reduction ($m=1$).
The next possible derivation expresses the fact that it is possible to reorder the elements to be reduced, expressing the commutativity property of our definition of reduction.
The third derivation is actually the only place where parallelism is expressed in the definition of our reduction pattern.
This rule expressed the fact that it is valid to partition the input elements first and then reduce them independently.
Finally, the last possible derivation expresses the notion that it is possible to perform a partial reduction with an iterative process by repetitively applying the same partial reduction function.
This concept is very important when considering how the reduction function is typically implemented on a GPU (iteratively reducing within a workgroup using the local memory).
%
\begin{rerule}{lcl}
  \pat{part-red}\ f\ z
    & \rightarrow &
      \pat{reduce}\ f\ z\\
    & | &
      \pat{part-red}\ f\ z \circ \pat{reorder}\\    
    & | &
      \pat{join}
        \circ \pat{map\ (part-red$\ f\ z$)}
        \circ \pat{split}\ n\\
    & | &
      \pat{iterate}\ n\ \pat{(part-red$\ f\ z$)}
  \label{eq:algo:part-red}
\end{rerule}


\paragraph{Simplification Rules}
\autoref{eq:algo:simpl} shows our simplification rules.
They express the fact that consecutive \pat{split}-\pat{join} pairs and \pat{asVector}-\pat{asScalar} pairs are equivalent to the identity. % function \emph{id}.
%
\begin{rerule}{lcl}
  \pat{split}\ n \circ \pat{join}\ n
    \hspace{1em} |\hspace{1em}
      \pat{join}\ n \circ \pat{split}\ n
        & \rightarrow & \pat{id}\\
  \pat{asVector}\ n \circ \pat{asScalar}\ n
    \hspace{1em} |\hspace{1em}
      \pat{asScalar}\ n \circ \pat{asVector}\ n
        & \rightarrow & \pat{id}
  \label{eq:algo:simpl}
\end{rerule}

\paragraph{Fusion Rules}
Finally, our fusion rules are shown in \autoref{eq:algo:fusion}.
The first rule fuses the functions applied by two consecutive maps.
The second rule fuses the map-reduce pattern by creating a lambda function that is the results of merging function $f$ and $g$ from the original reduction and map respectively.
This rule only applies to the sequential version since this is the only implementation not requiring the associativity property required by the more generic $\pat{reduce}$ pattern.
When generating code, these rules in effect allow us to fuse the implementation of the different functions and avoid having to store temporary results.
% More sophisticated rules exist in the functional programming community~\cite{fusion} which we want to incorporate in the future.
% The functional community has also studied rules for fusion (a.k.a. \emph{deforestation}) from a theoretical~\cite{} as well as a more practical point of view~\cite{jones01playing}.
The functional programming community has studied more generic rules for fusion~\cite{coutts07streamfusion,jones01playing}.
%More generic rules for fusion have been studied by the functional programming community~\cite{coutts07streamfusion,jones01playing}.
However, as we currently focus on a restricted set of patterns our simpler fusion rules have, so far, proven to be sufficient.
%
\begin{rerule}{lcl}
  \pat{map}\ f \circ \pat{map}\ g
    & \rightarrow & map\ (f \circ g)\\
  \pat{reduce-seq}\ f\ z \circ \pat{map-seq}\ g
    & \rightarrow & \\
  {\hspace{3em}}
  \pat{reduce-seq}\
    \big(\ \lambda\ acc,x\ .
      &\hspace{-.75em} f\ acc\ (g\ x)&\hspace{-.75em}\big)\ z
      % only reduce sequential is valid because non-associativity!!!
  \label{eq:algo:fusion}
\end{rerule}

\paragraph{Summary}
\todo{...}
\autoref{fig:algoRules} given an overview of all algorithmic rules.

\newlength{\ruleSpace}
\setlength{\ruleSpace}{0.5em}
\begin{figure}[t]
\centering
\begin{subfigure}[b]{1\linewidth}
  \begin{mdframed}
    \begin{rerule*}{lclcl}
          f & \rightarrow & f \circ \pat{id} & | & \pat{id} \circ f
    \end{rerule*}
  \end{mdframed}
  \vspace{-1em}
  \caption{Identity}
  \label{fig:algo:identity}
\end{subfigure}

\vspace{\ruleSpace}
\begin{subfigure}[b]{1\linewidth}
  \begin{mdframed}
    \begin{rerule*}{lcl}
      \pat{\textbf{iterate}}^{m+n}\ f
        & \rightarrow &
          \pat{\textbf{iterate}}^m\ f
            \circ \pat{\textbf{iterate}}^n\ f
    \end{rerule*}
  \end{mdframed}
  \vspace{-1em}
  \caption{Iterate decomposition}
  \label{fig:algo:iterate}
\end{subfigure}

\vspace{\ruleSpace}
\begin{subfigure}[b]{1\linewidth}
  \begin{mdframed}
    \begin{rerule*}{lcl}
      \pat{map}\ f \circ \pat{reorder}
        & \rightarrow & \pat{reorder} \circ \pat{map}\ f\\
      \pat{reorder} \circ \pat{map}\ f
        & \rightarrow & \pat{map}\ f \circ \pat{reorder}\\  
    \end{rerule*}
  \end{mdframed}
  \vspace{-1em}
  \caption{Reorder commutativity}
  \label{fig:algo:reorder}
\end{subfigure}

\vspace{\ruleSpace}
\begin{subfigure}[b]{1\linewidth}
  \begin{mdframed}
    \begin{rerule*}{lcl}
      \pat{map}\ f
        & \rightarrow &
          \pat{\textbf{join}}
            \circ \pat{map\ (map\ $f$)}
            \circ \pat{\textbf{split}}\ n
    \end{rerule*}
  \end{mdframed}
  \vspace{-1em}
  \caption{Split-join}
  \label{fig:algo:splitjoin}
\end{subfigure}

\vspace{\ruleSpace}
\begin{subfigure}[b]{1\linewidth}
  \begin{mdframed}
    \begin{rerule*}{lcl}
      \pat{reduce}\ f\ z
        & \rightarrow &
          \pat{reduce}\ f\ z \circ \pat{part-red}\ f\ z\\
      \pat{part-red}\ f\ z
        & \rightarrow &
          \pat{reduce}\ f\ z\\
        & | &
          \pat{part-red}\ f\ z \circ \pat{reorder}\\    
        & | &
          \pat{\textbf{join}}
            \circ \pat{map\ (part-red$\ f\ z$)}
            \circ \pat{\textbf{split}}\ n\\
        & | &
          \pat{\textbf{iterate}}\ n\ \pat{(part-red$\ f\ z$)}\\
    \end{rerule*}
  \end{mdframed}
  \vspace{-1em}
  \caption{Reduction}
  \label{fig:algo:red}
\end{subfigure}


\vspace{\ruleSpace}
\begin{subfigure}[b]{1\linewidth}
  \begin{mdframed}
    \begin{rerule*}{lcl}
      \pat{\textbf{split}}\ n \circ \pat{\textbf{join}}\ n
        \hspace{1em} |\hspace{1em}
          \pat{\textbf{join}}\ n \circ \pat{\textbf{split}}\ n
            & \rightarrow & \pat{id}\\
      \pat{\textbf{asVector}}\ n \circ \pat{\textbf{asScalar}}\ n
        \hspace{1em} |\hspace{1em}
          \pat{\textbf{asScalar}}\ n \circ \pat{\textbf{asVector}}\ n
            & \rightarrow & \pat{id}\\
    \end{rerule*}
  \end{mdframed}
  \vspace{-1em}
  \caption{Simplification rules}
  \label{fig:algo:simpl}
\end{subfigure}

\vspace{\ruleSpace}
\begin{subfigure}[b]{1\linewidth}
  \begin{mdframed}
    \begin{rerule*}{lcl}
      \pat{map}\ f \circ \pat{map}\ g
        & \rightarrow & map\ (f \circ g)\\
      \pat{\textbf{reduce-seq}}\ f\ z \circ \pat{\textbf{map-seq}}\ g
        & \rightarrow & \\
      {\hspace{3em}}
      \pat{\textbf{reduce-seq}}\
        \big(\ \lambda\ acc,x\ .
          &\hspace{-.75em} f\ acc\ (g\ x)&\hspace{-.75em}\big)\ z\\
          % only reduce sequential is valid because non-associativity!!!
    \end{rerule*}
  \end{mdframed}
  \vspace{-1em}
  \caption{Fusion rules}
  \label{fig:algo:fusion}
\end{subfigure}

\caption{Algorithmic rules. Bold patterns are known to the code generator.}
\label{fig:algoRules}
\end{figure}
\FloatBarrier





\subsection{OpenCL-Specific Rules}

In this section we discuss our \OpenCL-specific rules that are used to apply \OpenCL optimizations and to lower high-level algorithmic concepts down to \OpenCL-specific ones.
Patterns that are known to the code generator are shown in bold.% in both Figure~\ref{fig:algo} and~\ref{fig:low}.
%The code generation process is described separately in section~\ref{codegen}.

% \paragraph{Identity}
% Figure~\ref{fig:low:id} shows the rewrite rule for the identity pattern.
% We decided to only support code generation for the \pat{id-primitive} function
% which is valid only on non-array data types.
% % in our code generator which is the id function.
% We explicitly use the \pat{map} pattern in order to implement the identity function for arrays.

\paragraph{Maps}
The rule in \autoref{eq:low:map} is used to produce \OpenCL-specific map implementations that match the thread hierarchy of OpenCL.
Our implementation maintains context information to ensure the hierarchy is respected.
For instance, it is only legal to nest a \pat{map-local} inside a \pat{map-workgroup}.
%Similarly, a \pat{map-global} pattern can not be nested inside a \pat{map-workgroup} pattern as it represents global threads not organized in OpenCL workgroups.
%
% The \emph{pre} and \emph{post} conditions reflect the hierarchy presented in Figure~\ref{fig:map} corresponding to the OpenCL programming model.
% For instance for \pat{map-workgroup}, we have the pre condition that we must not be within a \pat{map-workgroup} ($\overline{mwg}$) already.
%
%Some rules can only be activated if certain conditions are true.
%We use the syntax $[pre:post]$ to represent the pre and post conditions of a rule.
%The pre condition $pre$ corresponds to the list of conditions that must be true for the rule to be applied.
%The post condition $post$ is set for any function bound to the pattern.
%The $\overline{\rule[-.3\baselineskip]{0pt}{1.5ex}\hspace{1.5ex}}$, $\wedge$, and $\vee$ corresponds to the logical \emph{not}, \emph{and}, and \emph{or} operators respectively.
%
\begin{rerule}{lclcl}
  \pat{map}\ f
    & \rightarrow & \pat{map-workgroup}\ f\\
    & | & \pat{map-local}\ f\\
    & | & \pat{map-global}\ f\\
    & | & \pat{map-seq}\ f
  \label{eq:low:map}
\end{rerule}

\paragraph{Reduction}
There is only one lowering rule for reduction (\autoref{eq:low:red}), which expresses the fact that the only implementation known to the code generator is a sequential reduction.
Parallel implementations are defined at a higher level by composition of other algorithmic patterns.
Most existing high performance compilers treat the reduction directly as an irreducible primitive operation.
%The power of our approach is that the code generator implementation only needs to know about the simple sequential reduction.
%As a result, it is possible to explore different implementation for the reduction by simply applying different rules.
With our approach it is possible to explore different implementations for the reduction by simply applying different rules.
%
\begin{rerule}{lcl}
  \pat{reduce}\ f\ z & \rightarrow & \pat{reduce-seq}\ f\ z
  \label{eq:low:red}
\end{rerule}


\paragraph{Reorder}
\autoref{eq:low:stride} presents the rule that reorders elements of an array.
In our current implementation, we support two types of reordering:
no reordering, represented by the \pat{id} function, and \pat{reorder-stride}, which reorders elements with a certain stride $s$.
As described earlier, the major use case for the stride reorder is to enable coalesced memory accesses.
%Note that other types of reordering functions could be implemented easily within this framework such as user-defined reorder function.
%
\begin{rerule}{lcl}
  \pat{reorder} & \rightarrow & \pat{reorder-stride}\ s\\
                & | & \pat{id}
  \label{eq:low:stride}
\end{rerule}

\paragraph{Local and Global Memory}
\ref{eq:low:mem} shows two rules that enable GPU local memory usage.
They express the fact that the result of a \pat{map-local} can always be stored in local memory or back in global memory.
This holds since a \pat{map-local} always exists within a \pat{map-workgroup} for which the local memory is defined.
These rules allow us to determine how the data is mapped to the GPU memory hierarchy.
%
\begin{rerule}{lcl}
  \pat{map-local}\ f & \rightarrow & \pat{toGlobal\ (map-local}\ f)\\  
  \pat{map-local}\ f & \rightarrow & \pat{toLocal\ (map-local}\ f)\\  
  \label{eq:low:mem}
\end{rerule}



\paragraph{Vectorization}
\autoref{eq:algo:vect} shows the vectorization rule.
Vectorization is achieved by using the \pat{asVector} and corresponding \pat{asScalar}, which changes the element type of an array and adjust the length accordingly.
This rule is only allowed to be applied once to a given \pat{map(f)} pattern.
This constrain can easily be checked by looking at the function's type. %; if it is a vector type, the rule cannot be applied.
Another set of rules, not shown here for space reason, is used to propagate the \pat{vect} function recursively within $f$.
%
\begin{rerule}{lcl}
  \pat{map}\ f
    & \rightarrow &
      \pat{asScalar}
        \circ \pat{map\ (vect}\ n\ f)
        \circ \pat{asVector}\ n
  \label{eq:algo:vect}
\end{rerule}

\paragraph{Summary}
\todo{...}
\autoref{fig:lowRules} show an overview of the \OpenCL-specific rules.

\begin{figure}[t]
\centering
\begin{subfigure}[b]{1\linewidth}
  \begin{mdframed}
    \begin{rerule*}{lclcl}
      \pat{map}\ f
        & \rightarrow &
          \pat{\textbf{map-workgroup}}\ f & | & \pat{\textbf{map-local}}\ f\\
        & | &
          \pat{\textbf{map-global}}\ f    & | & \pat{\textbf{map-seq}}\ f\\
    \end{rerule*}
  \end{mdframed}
  \vspace{-1em}
  \caption{Map}
  \label{fig:low:map}
\end{subfigure}

\vspace{\ruleSpace}
\begin{subfigure}[b]{1\linewidth}
  \begin{mdframed}
    \begin{rerule*}{lcl}
      \pat{reduce}\ f\ z
        & \rightarrow &
          \pat{\textbf{reduce-seq}}\ f\ z
    \end{rerule*}
  \end{mdframed}
  \vspace{-1em}
  \caption{Reduction}
  \label{fig:low:red}
\end{subfigure}

\vspace{\ruleSpace}
\begin{subfigure}[b]{1\linewidth}
  \begin{mdframed}
    \begin{rerule*}{lclcl}
      \pat{reorder}
        & \rightarrow &
          \pat{\textbf{reorder-stride}}\ s & | & \pat{id}
    \end{rerule*}
  \end{mdframed}
  \vspace{-1em}
  \caption{Stride accesses or normal accesses}
  \label{fig:low:stride}
\end{subfigure}

\vspace{\ruleSpace}
\begin{subfigure}[b]{1\linewidth}
  \begin{mdframed}
    \begin{rerule*}{lcl}
      \pat{\textbf{map-local}}\ f
        & \rightarrow &
          \pat{\textbf{toGlobal}\ (\textbf{map-local}}\ f)\\  
      \pat{\textbf{map-local}}\ f
        & \rightarrow & \pat{\textbf{toLocal}\ (\textbf{map-local}}\ f)\\  
    \end{rerule*}
  \end{mdframed}
  \vspace{-1em}
  \caption{Local/Global memory}
  \label{fig:low:mem}
\end{subfigure}

\vspace{\ruleSpace}
\begin{subfigure}[b]{1\linewidth}
  \begin{mdframed}
    \begin{rerule*}{lcl}
      \pat{map}\ f
        & \rightarrow &
          \pat{\textbf{asScalar}}
            \circ \pat{map\ (\textbf{vect}}\ n\ f)
            \circ \pat{\textbf{asVector}}\ n
    \end{rerule*}
  \end{mdframed}
  \vspace{-1em}
  \caption{Vectorization}
  \label{fig:algo:vect}
\end{subfigure}

\caption{OpenCL-specific rules. Bold patterns are known to the code generator.}
\label{fig:lowRules}
\end{figure}

\FloatBarrier


%\subsection{Automatic Rewriting Strategy}
%\label{sec:search}
%
%The rules presented in this section define a search space of possible implementations.
%In order to find the best possible low-level expression for a given target device, we have developed a simple automatic search strategy based loosely on Bandit-based optimization~\cite{demesmay09bandit}.
%Note that our current search strategy is just designed to prove that it is possible to find good implementations.
%We envision replacing this exploration strategy in the future by using machine-learning techniques to avoid having to search the space.
%However, this is orthogonal to the work presented in this paper.
%
%Our search strategy starts with the high-level expression and determines all the valid rules that can be applied at this stage.
%We use a Monte-Carlo method for evaluating the potential impact of each rule by randomly walking down the search tree.
%The rule that will lead to the best performance following the Monte-Carlo descent is chosen and applied to the expression.
%This process is repeated until we reach a terminal expression.
%Note that in addition to selecting the rules, we also search at the same time for the parameters controlling our patterns such as the vector size for the $\pat{vect}^n$ pattern.
%Using this simple strategy, we found that less than a thousand expressions were evaluated to reach a solution in most cases.

\subsection{Example: Deriving a Fused Implementation}
\label{sec:example}

\newcommand{\Reduce}{\text{\textit{reduce}}\xspace}
\newcommand{\PartRed}{\text{\textit{part-red}}\xspace}
\newcommand{\RedSeq}{\text{\textit{reduce-seq}}\xspace}
\newcommand{\Map}{\text{\textit{map}}\xspace}
\newcommand{\MapSeq}{\text{\textit{map-seq}}\xspace}
\newcommand{\MyJoin}{\text{\textit{join}}\xspace}
\newcommand{\MySplit}[1]{\text{\textit{split}}^{#1}\xspace}

\begin{figure*}[t]
\begin{align}
  \text{\textit{asum}}
  &\hspace{.2em}=\hspace{.2em} \Reduce\ (+)\ 0\ \circ\ \Map\ abs\\[1em]
  %
  & \overset{\ref{fig:algo:red}}{\hspace{.2em}=\hspace{.2em}}
      \begin{aligned}
        & \Reduce\ (+)\ 0\\
        &\circ\ \MyJoin\ \circ\ \Map\ (\PartRed\ (+)\ 0)\ \circ\ \MySplit\ n\\
        &\circ\ \Map\ abs
      \end{aligned}\\[1em]
  %
  & \overset{\ref{fig:algo:splitjoin}}{=\hspace{.2em}}
      \begin{aligned}
        & \Reduce\ (+)\ 0\\
        &\circ\ \MyJoin\ \circ\ \Map\ (\PartRed\ (+)\ 0)\ \circ\ \MySplit\ n\\
        &\circ\ \MyJoin\ \circ\ \Map\ (\Map\ abs)\ \circ\ \MySplit\ n
      \end{aligned}\\[1em]
  %
  & \overset{\ref{fig:algo:simpl}}{=\hspace{.2em}}
      \begin{aligned}
        & \Reduce\ (+)\ 0\\
        &\circ\ \MyJoin\ \circ\ \Map\ (\PartRed\ (+)\ 0)\\
        &\circ\ \Map\ (\Map\ abs)\ \circ\ \MySplit\ n
      \end{aligned}\\[1em]
  %
  & \overset{\ref{fig:algo:fusion}}{=\hspace{.2em}}
      \begin{aligned}
        & \Reduce\ (+)\ 0\\
        &\circ\ \MyJoin\ \circ\ \Map\ (\PartRed\ (+)\ 0\ \circ\ \Map\ abs)\\
        &\circ\ \MySplit\ n
      \end{aligned}\\[1em]
  %
  & \overset{\ref{fig:low:map}}{=\hspace{.2em}}
      \begin{aligned}
        & \Reduce\ (+)\ 0\\
        &\circ\ \MyJoin\ \circ\ \Map\ (\PartRed\ (+)\ 0\ \circ\ \MapSeq\ abs)\\
        &\circ\ \MySplit\ n
      \end{aligned}\\[1em]
  %
  & \overset{\hspace{-1.3em}\ref{fig:algo:red}\& \ref{fig:low:red}\hspace{-1.1em}}{=\hspace{.2em}}
      \begin{aligned}
        & \Reduce\ (+)\ 0\\
        &\circ\ \MyJoin\ \circ\ \Map\ (\RedSeq\ (+)\ 0\ \circ\ \MapSeq\ abs)\\
        &\circ\ \MySplit\ n
      \end{aligned}\\[1em]
  %
  & \overset{\ref{fig:algo:fusion}}{=\hspace{.2em}}
      \begin{aligned}
        & \Reduce\ (+)\ 0\\
        &\circ\ \MyJoin\\
        &\circ\ \Map\ \big(\RedSeq\ (\lambda\ a, b\ .\ a+(abs\ b))\ 0\big)\\
        &\circ\ \MySplit\ n
      \end{aligned}
\end{align}
\vspace{-2em}
\caption{Derivation for \emph{asum}$(\vec{x})$ to a fused parallel version.
  The numbers above the equality sign refer to the rules from sections ....
}
\label{fig:derivation}
\end{figure*}




To achieve good performance it is in general beneficial to avoid storing intermediate results.
Rule~\ref{fig:algo:fusion} allows us to apply this principle and fuse two patterns into one, thus, avoiding intermediate results.
Figure~\ref{fig:derivation} shows how we can derive a fused version for calculating the sum of absolute value, \emph{asum}, from the high-level expression written by the programmer.
We write the derivation as a sequence of equations using a slightly more mathematical notation.
The numbers above the equality sign refer to the rules applied.

We start by applying the reduction rule~\ref{fig:algo:red} twice:
first to replace \pat{reduce} with \pat{reduce}~$\circ$~\pat{part-red} and then a second time to expand \pat{par-red}.
To get (2) we expand \pat{map}, which can be simplified by removing the two corresponding \pat{join} and \pat{split} patterns.
In the step from (3) to (4) two \pat{map} patterns are fused and in the next step the nested \pat{map} is lowered into the \pat{map-seq} pattern to obtain (5).
By first transforming \pat{part-red} back into \pat{reduce} (using rule~\ref{fig:algo:red}) and then applying the lowering rule~\ref{fig:low:red} we get (6).
Finally, we apply rule~\ref{fig:algo:fusion} to fuse the \pat{map-seq} and \pat{reduce-seq} into a single \pat{reduce-seq}.
This sequence of transformations results in expression (7), which allows for a more optimal implementation since no temporary storage is required for the intermediate result.



\FloatBarrier



\subsection{Summary}

The power of our approach lies in the composition of our rules that produce complex low-level expressions from simple high-level expressions.
Looking back at our example in Figure~\ref{fig:codeex}, we see how a simple algorithmic pattern can effectively be derived into a low-level expression by applying the rules.
This expression matches hardware concepts expressible with OpenCL such as mapping computation and data to the thread and memory hierarchy. % and vectorization.
Each single rule encodes a simple, easy to understand, provable fact.
By composition of the rules we systematically derive low-level expressions which are semantically equivalent to the high-level expressions by construction.
This results in a powerful mechanism to safely explore the space of possible implementations.
% Using these rules and our automatic search technique we now have a powerful mechanism to explore the space of possible program implementations.

% \from{PACT begin}
% \paragraph{Rewrite Rules (PACT)}
% 
% This section introduces our set of rewrite rules that transform high-level expressions written using our algorithmic patterns into semantically equivalent expressions.
% One goal of our approach is to keep each rule as simple as possible and only express one fundamental concept at a time.
% For instance the vectorization rule, as we will see, is the only place where we express the vectorization concept.
% This is different from most prior approaches that would for instance produce a special vectorized version of different algorithmic patterns such as map or reduce.
% The superiority of using such an approach lies in the power of composition;
% many rules can be applied successively to produce expressions that compose hardware concepts or optimizations and that are correct by construction.
% 
% Similarly to our patterns, we distinguish between algorithmic and lowering rules.
% Algorithmic rules produces derivations that represent the different algorithmic choices and are shown in Figure~\ref{fig:algo}.
% The lowering rules shown in Figure~\ref{fig:low} map expressions to hardware patterns expressible with the OpenCL programming model.
% Once the expression is in its lowest form, it is possible to produce OpenCL code easily with our code generator as seen in the previous section.
% 
% 
% \paragraph{Syntax and Rule Derivation}
% Some rules can only be activated if certain conditions are true.
% The syntax $[pre:post]$ represents the pre and post conditions.
% The pre condition $pre$ corresponds to the list of conditions that must be true for the rule to be applied.
% The post condition $post$ is set for any function bound to the pattern.
% The $\overline{\rule[-.3\baselineskip]{0pt}{1.5ex}\hspace{1.5ex}}$, $\wedge$, and $\vee$ corresponds to the logical \emph{not}, \emph{and}, and \emph{or} operators respectively.
% 
% We use leftmost derivation when applying the rules, which means that the leftmost non-terminal is always derived first.
% 
% 
% \subsection{Algorithmic Rules}
% 
% \begin{figure}[t]
% \centering
% \begin{subfigure}[b]{1\linewidth}
% \begin{mdframed}
% $$
% \begin{array}{lllll}
%   f & \rightarrow & f \circ \pat{id} & | & \pat{id} \circ f\\
% \end{array}
% $$
% \end{mdframed}
%   \caption{Identity}
%   \label{fig:algo:identity}
% \end{subfigure}
% 
% \vspace{-0.5em}
% \begin{subfigure}[b]{1\linewidth}
% \begin{mdframed}
% $$
% \begin{array}{lrl}
%   \pat{\textbf{iterate}}^{m+n}\pat{(f)} & \rightarrow & \pat{\textbf{iterate}}^m\pat{(f)} \circ \pat{\textbf{iterate}}^n\pat{(f)}\\
%   \end{array}
% $$
% \end{mdframed}
%   \caption{Iterate decomposition}
%   \label{fig:algo:iterate}
% \end{subfigure}
% 
% \vspace{-0.5em}
% \begin{subfigure}[b]{1\linewidth}
% \begin{mdframed}
% $$
% \begin{array}{lrl}
%   \pat{map(f)} \circ \pat{reorder} & \rightarrow & \pat{reorder} \circ \pat{map(f)}\\
%   \pat{reorder} \circ \pat{map(f)} & \rightarrow & \pat{map(f)} \circ \pat{reorder}\\  
% \end{array}
% $$
% \end{mdframed}
%   \caption{Reorder commutativity}
%   \label{fig:algo:reorder}
% \end{subfigure}
% 
% \vspace{-0.5em}
% \begin{subfigure}[b]{1\linewidth}
% \begin{mdframed}
% $$
% \begin{array}{lrl}
%   \pat{map(f)} & \rightarrow & \pat{\textbf{outerJoin}} \circ \pat{map(map(f))}\circ \pat{\textbf{outerSplit}}^n
% \end{array}
% $$
% \end{mdframed}
%   \caption{Split-join}
%   \label{fig:algo:splitjoin}
% \end{subfigure}
% 
% \vspace{-0.5em}
% \begin{subfigure}[b]{1\linewidth}
% \begin{mdframed}
% $$
% \begin{array}{lrl}
% \multicolumn{3}{l}{\pat{part-red(f,z)} \rightarrow} \\
%   & & \pat{reduce(f,z)}\\                                                      
%   &| & \pat{part-red(f,z)} \circ \pat{reorder}\\    
%   &| & \pat{\textbf{outerJoin}} \circ \pat{map(part-red(f,z))} \circ \pat{\textbf{outerSplit}}^n\\
%   &| & \pat{\textbf{iterate}}^{n}\pat{(part-red(f,z))}\\
% \end{array}
% $$
% $$
% \begin{array}{lrl}
%   \pat{reduce(f,z)} & \rightarrow & \pat{reduce(f,z)} \circ \pat{part-red(f,z)}\\
%   & | & \pat{\textbf{iterate}}^{\infty}\pat{(part-red(f,z))}\\
% 
% \end{array}
% $$
% \end{mdframed}
%   \caption{Reduction}
%   \label{fig:algo:red}
% \end{subfigure}
% 
% \vspace{-0.5em}
% \begin{subfigure}[b]{1\linewidth}
% \begin{mdframed}
% $$
% \begin{array}{lrll}
% \multicolumn{2}{l}{\pat{map(f)} \rightarrow} \\
%  {\ }   & \pat{\textbf{innerJoin}} \circ \pat{map(\textbf{vect}}^n\pat{(f))} \circ \pat{\textbf{innerSplit}}^n & [\hspace{.2em}\overline{vec}:vec\hspace{.2em}]
% \end{array}
% $$
% \end{mdframed}
%   \caption{Vectorization}
%    \label{fig:algo:vect}
% \end{subfigure}
% 
% \vspace{-0.5em}
% \begin{subfigure}[b]{1\linewidth}
% \begin{mdframed}
% $$
% \begin{array}{lrl}
% f \circ \pat{id}\ |\ \pat{id} \circ f & \rightarrow & f\\
% \pat{\textbf{outerSplit}}^n \circ \pat{\textbf{outerJoin}}^n & \rightarrow & \pat{id}\\
% \pat{\textbf{innerSplit}}^n \circ \pat{\textbf{innerJoin}}^n & \rightarrow & \pat{id}\\
% \end{array}
% $$
% \end{mdframed}
%   \caption{Simplification rules}
%    \label{fig:algo:simpl}
% \end{subfigure}
% 
% \vspace{-0.5em}
% \begin{subfigure}[b]{1\linewidth}
% \begin{mdframed}
% $$
% \begin{array}{lll}
% \pat{map(f)} \circ \pat{map(g)}                                & \rightarrow & map(f \circ g)\\
% \pat{\textbf{reduce-seq}(f,z)} \circ \pat{\textbf{map-seq}(g)} & \rightarrow & \\
% {\ \ } \pat{\textbf{reduce-seq}}(\lambda\ acc,x: f(acc,g(x)), z)\\
% \end{array}
% $$
% \end{mdframed}
%   \caption{Fusion rules}
%    \label{fig:algo:fusion}
% \end{subfigure}
% \vspace{-2em}
% \caption{Algorithmic rules. Bold patterns are known to the code generator.}
% \label{fig:algo}
% \end{figure}
% 
% \begin{figure}[t]
% \centering
% 
% \begin{subfigure}[b]{1\linewidth}
% \begin{mdframed}
% $$
% \begin{array}{lll}
% \multicolumn{3}{l}{\pat{map(f)} \rightarrow} \\
%        & \pat{\textbf{map-workgroup}(f)} & [\hspace{.2em}\overline{mwg}:mwg\hspace{.2em}]\\
% {\ } | & \pat{\textbf{map-local}(f)    } & [\hspace{.2em}\overline{mlc}\hspace{.7em} \wedge mwg\wedge\overline{mwp}:mlc\hspace{.2em}]\\
% {\ } | & \pat{\textbf{map-warp}(f)     } & [\hspace{.2em}\overline{mwp}\hspace{.2em} \wedge mwg\wedge\overline{mlc}:mwp\hspace{.2em}]\\
% {\ } | & \pat{\textbf{map-lane}(f)     } & [\hspace{.2em}\overline{mln}\hspace{.5em} \wedge mwp:mln\hspace{.2em}]\\
% {\ } | & \pat{\textbf{map-seq}(f)   } & [\hspace{.2em}\overline{msr}\hspace{.5em} \wedge (mlc\vee mln):msr\hspace{.2em}]\\          
% \end{array}
% $$
% \end{mdframed}
%   \caption{Map}
%   \label{fig:low:map}
% \end{subfigure}
% 
% \vspace{-0.5em}
% \begin{subfigure}[b]{1\linewidth}
% \begin{mdframed}
% $$
% \begin{array}{llll}
%   \pat{reduce(f,z)} & \rightarrow & \pat{\textbf{reduce-seq}(f,z)} & [\hspace{.2em}mlc\vee mln:\ \hspace{.2em}]\\   
% \end{array}
% $$
% \end{mdframed}
%   \caption{Reduction}
%   \label{fig:low:red}
% \end{subfigure}
% 
% \vspace{-0.5em}
% \begin{subfigure}[b]{1\linewidth}
% \begin{mdframed}
% $$
% \begin{array}{lrl}
%   \pat{\textbf{map-local}(f)} & \rightarrow & \pat{\textbf{asGlobal}(map-local(f))}\\  
%   \pat{\textbf{map-local}(f)} & \rightarrow & \pat{\textbf{asLocal}(map-local(f))}\\  
% \end{array}
% $$
% \end{mdframed}
%   \caption{Local/Global memory}
%   \label{fig:low:mem}
% \end{subfigure}
% 
% \vspace{-0.5em}
% \begin{subfigure}[b]{1\linewidth}
% \begin{mdframed}
% $$
% \begin{array}{lllll}
%   \pat{reorder}  & \rightarrow & \pat{\textbf{reorder-stride}}^s & | & \pat{id}
% \end{array}
% $$
% \end{mdframed}
%   \caption{Stride accesses or normal accesses}
%   \label{fig:low:stride}
% \end{subfigure}
% \vspace{-2em}
% \caption{Lowering rules. Bold patterns are known to the code generator.}
% \label{fig:low}
% \end{figure}
% 
% 
% 
% The first three rules (see Figure~\ref{fig:algo:identity},~\ref{fig:algo:iterate} and~\ref{fig:algo:reorder}) express the properties of the identity, iterate and reorder pattern.
% 
% 
% \paragraph{Split-join}
% The split-join rule in Figure~\ref{fig:algo:splitjoin} partitions a map into two maps.
% This allows us to nest maps in each other and, thus, \emph{map} the computation to the thread hierarchy of the OpenCL programming model such as \pat{map-workgroup(map-local(f))} as seen in our motivation example (Figure~\ref{fig:codeex}).
% 
% 
% \paragraph{Reduction}
% The reduction (and associated partial reduction) in Figure~\ref{fig:algo:red} is currently our most complex rule but also the most powerful one.
% It allows us to express the reduction function as a composition of other primitive functions, which is a fundamental aspect of our work.
% From the algorithmic point of view we first define a partial reduction pattern \pat{part-red}.
% This partial reduction reduces an array of $n$ elements to an array of $m$ elements where $1 \leq m < n$.
% The reduction can be derived in a partial reduction combined with a full reduction which ensures we end up with one unique element.
% Another possible derivation consists in iterating a partial reduction until a full reduction is achieved (this is what $\infty$ represents).
% 
% Note that our definition of \pat{reduce} remains correct since the result of a partial reduction is always composed with the reduction to ensure we end up with one unique element.
% 
% \paragraph{Partial Reduction}
% The first possible derivation for partial reduction, in Figure~\ref{fig:algo:red}, leads to the full reduction which means $m=1$.
% The next possible derivation expresses the fact that it is possible to reorder the elements to be reduced, expressing the commutativity property of our definition of reduction.
% The fourth derivation is actually the only place where parallelism is expressed in the definition of our reduction pattern.
% This rule expressed the fact that it is valid to partition the input elements first and then reduce them independently.
% Finally, the last possible derivation expresses the notion that it is possible to perform a partial reduction with an iterative process by repetitively applying the same partial reduction function.
% This concept is very important when considering how the reduction function is typically implemented on a GPU (iteratively reducing within a workgroup using the local memory).
% 
% \paragraph{Vectorization}
% Figure~\ref{fig:algo:vect} shows the vectorization rule.
% Vectorization is achieved by using the \pat{innerSplit} and corresponding \pat{innerJoin} which split the data on the innermost dimension only.
% The rule \emph{pre} and \emph{post} conditions ensure that vectorization is applied only once in any expression derived from $f$.
% Another set of rules, not shown here for space reason, are used to propagate the $vect^n$ function recursively within $f$.
% 
% \paragraph{Simplification Rules}
% Figure~\ref{fig:algo:simpl} shows our simplification rules.
% They express the property of the $id$ function and the removal of consecutive splits and joins.
% 
% \paragraph{Fusion Rules}
% Finally, our fusion rules are shown in Figure~\ref{fig:algo:fusion}.
% The first rule allows us to fuse the functions applied by two consecutive maps.
% The second rule fuses the map-reduce pattern by creating a lambda function that is the results of merging function $f$ and $g$ from the original reduction and map respectively.
% Note that this rule only applies to the sequential version since this is the only implementation that does not require the associativity property required by the more generic $\pat{reduce}$ pattern.
% When generating code, these rules in effect allow us to fuse the implementation of the different functions and avoid having to store temporary results.
% 
% 
% 
% 
% 
% \subsection{Lowering Rules}
% 
% Our lowering rules are shown in Figure~\ref{fig:low}.
% The purpose of these rules is to lower high-level expressions to low-level expressions, which consist only of hardware patterns for which we can generate code.
% These hardware patterns are highlighted in bold in Figure~\ref{fig:algo} and~\ref{fig:low}.
% 
% \paragraph{Maps}
% The rule shown in Figure~\ref{fig:low:map} is used to produce the different low-level \pat{maps}.
% The \emph{pre} and \emph{post} conditions reflect the hierarchy presented in Figure~\ref{fig:map} corresponding to the OpenCL programming model.
% For instance for \pat{map-workgroup}, we have the pre condition that we must not be within a \pat{map-workgroup} ($\overline{mwg}$) already.
% 
% \paragraph{Reduction}
% There is only one lowering rule for reduction (Figure~\ref{fig:low:red}), which expresses the fact that the only OpenCL implementation known to the code generator is a sequential reduction.
% The reduction pattern is defined at a higher level by composing other algorithmic patterns.
% To the best of our knowledge, all other existing high performance compilers treat the reduction directly as an irreducible primitive operation.
% The power of our approach is that the code generator implementation only needs to know about the simple sequential reduction.
% As a result, it is possible to explore different implementation for the reduction automatically by simply applying different rules.
% 
% \paragraph{Local/Global}
% Figure~\ref{fig:low:mem} shows the two rules that allow us to make use of the GPU local memory.
% These rules express the fact that the result of a \pat{map-local} can always be stored in local memory or back in global memory.
% This holds because a \pat{map-local} always exists within a \pat{map-workgroup} where the local memory is defined.
% These rules allow us to determine how the data is mapped to the GPU memory hierarchy.
% 
% \paragraph{Reorder}
% Finally, Figure~\ref{fig:low:stride} presents the rule that reorders elements of an array.
% In our current implementation, we support two types of reordering:
% no reordering, represented by the \pat{id} function, and \pat{reorder-stride} which reorders elements with a certain stride $s$.
% As described earlier, the major use case for the stride reorder is to enable coalesced memory accesses.
% 
% 
% \subsection{Automatic Rewriting Strategy}
% \label{sec:search}
% 
% The rules presented in this section define a search space of possible implementations.
% In order to find the best possible low-level expression for a given target device, we have developed a simple automatic search strategy based loosely on Bandit-based optimization~\cite{demesmay09bandit}.
% Note that our current search strategy is just designed to prove that it is possible to find good implementations.
% We envision replacing this exploration strategy in the future by using machine-learning techniques to avoid having to search the space.
% However, this is orthogonal to the work presented in this paper.
% 
% Our search strategy starts with the high-level expression and determines all the valid rules that can be applied at this stage.
% We use a Monte-Carlo method for evaluating the potential impact of each rule by randomly walking down the search tree.
% The rule that will lead to the best performance following the Monte-Carlo descent is chosen and applied to the expression.
% This process is repeated until we reach a terminal expression.
% Note that in addition to selecting the rules, we also search at the same time for the parameters controlling our patterns such as the vector size for the $\pat{vect}^n$ pattern.
% Using this simple strategy, we found that less than a thousand expressions were evaluated to reach a solution in most cases.
% 
% 
% \subsection{Summary}
% 
% The power of our approach lies in the composition of our rules that produce complex expressions.
% Looking back at our motivation example in Figure~\ref{fig:codeex}, we can see how a simple algorithmic pattern such as \pat{map} can be effectively derived into a low-level expression by applying our rules.
% This expression matches various hardware concepts expressible with the OpenCL programming model such as mapping computation and data to the GPU thread and memory hierarchy respectively and vectorization.
% Using these rules and our automatic search technique we now have a powerful mechanism to explore the space of possible program implementations.
% \from{PACT end}
% 
% \from{PPoPP begin}
% \paragraph{Rewrite Rules (PPoPP)}
% This section introduces our set of rewrite rules that transform high-level expressions written using our algorithmic patterns into semantically equivalent expressions.
% One goal of our approach is to keep each rule as simple as possible and only express one fundamental concept at a time.
% For instance the vectorization rule, as we will see, is the only place where we express the vectorization concept.
% This is different from most prior approaches that would produce a special vectorized version of different algorithmic patterns such as map or reduce.
% The superiority of our approach lies in the power of composition;
% many rules can be applied successively to produce expressions that compose hardware concepts or optimizations and that are provably correct by construction.
% 
% Similarly to our patterns, we distinguish between algorithmic and lowering rules.
% Algorithmic rules produces derivations that represent the different algorithmic choices and are shown in Figure~\ref{fig:algo}.
% %The lowering rules map expressions to hardware patterns.
% Figure~\ref{fig:low} shows our OpenCL-specific rules which map expressions to OpenCL patterns.
% %are expressible with the OpenCL programming model.
% Once the expression is in its lowest form, it is possible to produce OpenCL code for each single pattern easily with our code generator as described in the previous section.
% 
% 
% %\paragraph{Syntax and Rule Derivation}
% %Some rules can only be activated if certain conditions are true.
% %We use the syntax $[pre:post]$ to represent the pre and post conditions of a rule.
% %The pre condition $pre$ corresponds to the list of conditions that must be true for the rule to be applied.
% %The post condition $post$ is set for any function bound to the pattern.
% %The $\overline{\rule[-.3\baselineskip]{0pt}{1.5ex}\hspace{1.5ex}}$, $\wedge$, and $\vee$ corresponds to the logical \emph{not}, \emph{and}, and \emph{or} operators respectively.
% 
% %We use leftmost derivation when applying the rules, which means that the leftmost non-terminal is always derived first.
% 
% 
% \subsection{Algorithmic Rules}
% 
% 
% 
% % The first three rules (see Figure~\ref{fig:algo:identity},~\ref{fig:algo:iterate} and~\ref{fig:algo:reorder}) express the properties of the identity, iterate and reorder pattern.
% 
% % \paragraph{Identity}
% % The identity rule in Figure~\ref{fig:algo:identity} specifies that it is always valid to compose any function with the identity function \pat{id}.
% % The \pat{id} function can acts as a copy operation; this is useful to express copy to local memory for instance when composed with the corresponding lowering rule.
%  
% \paragraph{Iterate decomposition}
% The rule in Figure~\ref{fig:algo:iterate} expresses the fact an iteration can be decomposed into several iterations.
% 
% \paragraph{Reorder commutativity}
% Figure~\ref{fig:algo:reorder} shows a rule stating that if the data can be reordered arbitrarily it does not matter if we apply a function $f$ to each element before or after the reordering.
% 
% \paragraph{Split-join}
% The split-join rule in Figure~\ref{fig:algo:splitjoin} partitions a map into two maps.
% This allows us to nest map patterns in each other and, thus, \emph{map} the computation to the thread hierarchy of the OpenCL programming model such as \pat{map-workgroup(map-local(f))} as seen in our motivation example (Figure~\ref{fig:codeex}).
% 
% 
% \paragraph{Reduction}
% The reduction (and associated partial reduction) in Figure~\ref{fig:algo:red} is currently our most complex rule but also the most powerful one.
% It expresses the reduction function as a composition of other primitive functions, which is a fundamental aspect of our work.
% From the algorithmic point of view we first define a partial reduction pattern \pat{part-red}.
% This partial reduction reduces an array of $n$ elements to an array of $m$ elements where $1 \leq m < n$.
% The reduction can be derived in a partial reduction combined with a full reduction which ensures we end up with one unique element.
% % Another possible derivation consists in iterating a partial reduction until a full reduction is achieved (this is what $\infty$ represents).
% %Note that our definition of \pat{reduce} remains correct since the result of a partial reduction is always composed with the reduction to ensure we end up with one unique element.
% 
% \paragraph{Partial Reduction}
% The first possible derivation for partial reduction, in Figure~\ref{fig:algo:red}, leads to the full reduction which means $m=1$.
% %The first two possible derivations for partial reduction, in figure~\ref{fig:algo:red}, lead to the $id$ function or the full reduction which means $m=n$ or $m=1$ respectively.
% The next possible derivation expresses the fact that it is possible to reorder the elements to be reduced, expressing the commutativity property of our definition of reduction.
% The third derivation is actually the only place where parallelism is expressed in the definition of our reduction pattern.
% This rule expressed the fact that it is valid to partition the input elements first and then reduce them independently.
% Finally, the last possible derivation expresses the notion that it is possible to perform a partial reduction with an iterative process by repetitively applying the same partial reduction function.
% This concept is very important when considering how the reduction function is typically implemented on a GPU (iteratively reducing within a workgroup using the local memory).
% 
% 
% \paragraph{Simplification Rules}
% Figure~\ref{fig:algo:simpl} shows our simplification rules.
% They express the fact that consecutive \pat{split}-\pat{join} pairs and \pat{asVector}-\pat{asScalar} pairs are equivalent to the identity function \emph{id}.
% 
% \paragraph{Fusion Rules}
% Finally, our fusion rules are shown in Figure~\ref{fig:algo:fusion}.
% The first rule fuses the functions applied by two consecutive maps.
% The second rule fuses the map-reduce pattern by creating a lambda function that is the results of merging function $f$ and $g$ from the original reduction and map respectively.
% This rule only applies to the sequential version since this is the only implementation not requiring the associativity property required by the more generic $\pat{reduce}$ pattern.
% When generating code, these rules in effect allow us to fuse the implementation of the different functions and avoid having to store temporary results.
% % More sophisticated rules exist in the functional programming community~\cite{fusion} which we want to incorporate in the future.
% % The functional community has also studied rules for fusion (a.k.a. \emph{deforestation}) from a theoretical~\cite{} as well as a more practical point of view~\cite{jones01playing}.
% More generic rules for fusion have been studies by the functional programming community~\cite{coutts07streamfusion,jones01playing}.
% However, as we currently focus on a restricted set of patterns our simpler fusion rules have, so far, proven to be sufficient.
% 
% 
% 
% \newlength{\ruleSpace}
% \setlength{\ruleSpace}{-1em}
% \begin{figure}[t]
% \centering
% % \begin{subfigure}[b]{1\linewidth}
% % \begin{mdframed}
% % $$
% % \begin{array}{lllll}
% %   f & \rightarrow & f \circ \pat{id} & | & \pat{id} \circ f\\
% % \end{array}
% % $$
% % \end{mdframed}
% %   \caption{Identity}
% %   \label{fig:algo:identity}
% % \end{subfigure}
% 
% %\vspace{\ruleSpace}
% \begin{subfigure}[b]{1\linewidth}
% \begin{mdframed}
% $$
% \begin{array}{lrl}
%   \pat{\textbf{iterate}}^{m+n}\pat{(f)} & \rightarrow & \pat{\textbf{iterate}}^m\pat{(f)} \circ \pat{\textbf{iterate}}^n\pat{(f)}\\
%   \end{array}
% $$
% \end{mdframed}
%   \caption{Iterate decomposition}
%   \label{fig:algo:iterate}
% \end{subfigure}
% 
% \vspace{\ruleSpace}
% \begin{subfigure}[b]{1\linewidth}
% \begin{mdframed}
% $$
% \begin{array}{lrl}
%   \pat{map(f)} \circ \pat{reorder} & \rightarrow & \pat{reorder} \circ \pat{map(f)}\\
%   \pat{reorder} \circ \pat{map(f)} & \rightarrow & \pat{map(f)} \circ \pat{reorder}\\  
% \end{array}
% $$
% \end{mdframed}
%   \caption{Reorder commutativity}
%   \label{fig:algo:reorder}
% \end{subfigure}
% 
% \vspace{\ruleSpace}
% \begin{subfigure}[b]{1\linewidth}
% \begin{mdframed}
% $$
% \begin{array}{lrl}
%   \pat{map(f)} & \rightarrow & \pat{\textbf{join}} \circ \pat{map(map(f))}\circ \pat{\textbf{split}}^n
% \end{array}
% $$
% \end{mdframed}
%   \caption{Split-join}
%   \label{fig:algo:splitjoin}
% \end{subfigure}
% 
% \vspace{\ruleSpace}
% \begin{subfigure}[b]{1\linewidth}
% \begin{mdframed}
% $$
% \begin{array}{lrl}
% \pat{reduce(f,z)} & \rightarrow & \pat{reduce(f,z)} \circ \pat{part-red(f,z)}\\[.5em]
% \pat{part-red(f,z)} & \rightarrow & \pat{reduce(f,z)}\\                                                      
%   &| & \pat{part-red(f,z)} \circ \pat{reorder}\\    
%   &| & \pat{\textbf{join}} \circ \pat{map(part-red(f,z))} \circ \pat{\textbf{split}}^n\\
%   &| & \pat{\textbf{iterate}}^{n}\pat{(part-red(f,z))}\\
% \end{array}
% %$$
% %$$
% %\begin{array}{lrl}
% %  \pat{reduce(f,z)} & \rightarrow & \pat{reduce(f,z)} \circ \pat{part-red(f,z)}%\\
% %  %& | & \pat{\textbf{iterate}}^{\infty}\pat{(part-red(f,z))}\\
% %
% %\end{array}
% $$
% \end{mdframed}
%   \caption{Reduction}
%   \label{fig:algo:red}
% \end{subfigure}
% 
% 
% \vspace{\ruleSpace}
% \begin{subfigure}[b]{1\linewidth}
% \begin{mdframed}
% $$
% \begin{array}{lrl}
% %f \circ \pat{id}\hspace{1em} |\hspace{1em} \pat{id} \circ f & \rightarrow & f\\
% \pat{\textbf{split}}^n \circ \pat{\textbf{join}}^n \hspace{1em} |\hspace{1em} \pat{\textbf{join}}^n \circ \pat{\textbf{split}}^n & \rightarrow & \pat{id}\\
% \pat{\textbf{asVector}}^n \circ \pat{\textbf{asScalar}}^n \hspace{1em} |\hspace{1em} \pat{\textbf{asScalar}}^n \circ \pat{\textbf{asVector}}^n       & \rightarrow & \pat{id}\\
% \end{array}
% $$
% \end{mdframed}
%   \caption{Simplification rules}
%    \label{fig:algo:simpl}
% \end{subfigure}
% 
% \vspace{\ruleSpace}
% \begin{subfigure}[b]{1\linewidth}
% \begin{mdframed}
% $$
% \begin{array}{lll}
% \pat{map(f)} \circ \pat{map(g)}                                & \rightarrow & map(f \circ g)\\
% \pat{\textbf{reduce-seq}(f,z)} \circ \pat{\textbf{map-seq}(g)} & \rightarrow & \\
% {\hspace{2em}} \pat{\textbf{reduce-seq}}(\lambda\ acc,x: f(acc,g(x)), z)\\  % only reduce sequential is valid because non-associativity!!!
% \end{array}
% $$
% \end{mdframed}
%   \caption{Fusion rules}
%    \label{fig:algo:fusion}
% \end{subfigure}
% \vspace{-2em}
% \caption{Algorithmic rules. Bold patterns are known to the code generator.}
% \label{fig:algo}
% \end{figure}
% 
% 
% 
% \begin{figure}[t]
% \centering
% 
% % \begin{subfigure}[b]{1\linewidth}
% % \begin{mdframed}
% % $$
% % \begin{array}{llll}
% %  \langle primitive \rangle & \pat{id} & \rightarrow & \pat{\textbf{id-primitive}}\\
% %  \langle T[\ ] \rangle     & \pat{id} & \rightarrow & \pat{map(id)}\\  
% % \end{array}
% % $$
% % \end{mdframed}
% %   \caption{Identity}
% %   \label{fig:low:id}
% % \end{subfigure}
% 
% \begin{subfigure}[b]{1\linewidth}
% \begin{mdframed}
% $$
% %\begin{array}{lll}
% %\multicolumn{3}{l}{\pat{map(f)} \rightarrow} \\
% %       & \pat{\textbf{map-workgroup}(f)} & [\hspace{.2em}\overline{mwg}\hspace{.2em}\wedge\hspace{.2em}\overline{mgl}:mwg\hspace{.2em}]\\
% %{\ } | & \pat{\textbf{map-local}(f)    } & [\hspace{.2em}\overline{mlc} \hspace{.2em}\wedge\hspace{.2em} mwg \hspace{.2em}\wedge\hspace{.2em} \overline{mgl}: mlc\hspace{.2em}]\\
% %%{\ } | & \pat{\textbf{map-warp}(f)     } & [\hspace{.2em}\overline{mwp}\hspace{.2em} \wedge mwg\wedge\overline{mlc}:mwp\hspace{.2em}]\\
% %%{\ } | & \pat{\textbf{map-lane}(f)     } & [\hspace{.2em}\overline{mln}\hspace{.5em} \wedge mwp:mln\hspace{.2em}]\\
% %{\ } | & \pat{\textbf{map-global}(f)   } & [\hspace{.2em}\overline{mgl} \hspace{.2em}\wedge\hspace{.2em} \overline{mlc}\hspace{.2em}\wedge\hspace{.2em}\overline{mwg}:mgl\hspace{.2em}]\\
% %{\ } | & \pat{\textbf{map-seq}(f)      } & [\hspace{.2em}\overline{msq}\hspace{.5em} \wedge (mlc\vee mln):msq\hspace{.2em}]\\          
% %\end{array}
% \begin{array}{lllll}
% \pat{map(f)} & \rightarrow & \pat{\textbf{map-workgroup}(f)} & | & \pat{\textbf{map-local}(f)}\\
%  & | & \pat{\textbf{map-global}(f)} & | & \pat{\textbf{map-seq}(f)}\\          
% \end{array}
% $$
% \end{mdframed}
%   \caption{Map}
%   \label{fig:low:map}
% \end{subfigure}
% 
% \vspace{\ruleSpace}
% \begin{subfigure}[b]{1\linewidth}
% \begin{mdframed}
% $$
% \begin{array}{lll}
%   \pat{reduce(f,z)} & \rightarrow & \pat{\textbf{reduce-seq}(f,z)}% & [\hspace{.2em}mlc\vee mln:\ \hspace{.2em}]\\   
% \end{array}
% $$
% \end{mdframed}
%   \caption{Reduction}
%   \label{fig:low:red}
% \end{subfigure}
% 
% \vspace{\ruleSpace}
% \begin{subfigure}[b]{1\linewidth}
% \begin{mdframed}
% $$
% \begin{array}{lllll}
%   \pat{reorder}  & \rightarrow & \pat{\textbf{reorder-stride}}^s & | & \pat{id}
% \end{array}
% $$
% \end{mdframed}
%   \caption{Stride accesses or normal accesses}
%   \label{fig:low:stride}
% \end{subfigure}
% 
% \vspace{\ruleSpace}
% \begin{subfigure}[b]{1\linewidth}
% \begin{mdframed}
% $$
% \begin{array}{lrl}
%   \pat{\textbf{map-local}(f)} & \rightarrow & \pat{\textbf{toGlobal}(\textbf{map-local}(f))}\\  
%   \pat{\textbf{map-local}(f)} & \rightarrow & \pat{\textbf{toLocal}(\textbf{map-local}(f))}\\  
% \end{array}
% $$
% \end{mdframed}
%   \caption{Local/Global memory}
%   \label{fig:low:mem}
% \end{subfigure}
% 
% \vspace{\ruleSpace}
% \begin{subfigure}[b]{1\linewidth}
% \begin{mdframed}
% $$
% %\begin{array}{lrll}
% %\multicolumn{2}{l}{\pat{map(f)} \rightarrow} \\
% % {\ }   & \pat{\textbf{asScalar}} \circ \pat{map(\textbf{vect}}^n\pat{(f))} \circ \pat{\textbf{asVector}}^n & [\hspace{.2em}\overline{vec}:vec\hspace{.2em}]
% %\end{array}
% \begin{array}{lrll}
% \pat{map(f)} & \rightarrow & \pat{\textbf{asScalar}} \circ \pat{map(\textbf{vect}}^n\pat{(f))} \circ \pat{\textbf{asVector}}^n
% \end{array}
% $$
% \end{mdframed}
%   \caption{Vectorization}
%    \label{fig:algo:vect}
% \end{subfigure}
% 
% \vspace{-2em}
% \caption{OpenCL-specific rules. Bold patterns are known to the code generator.}
% \label{fig:low}
% \end{figure}
% 
% 
% 
% 
% 
% \subsection{OpenCL-Specific Rules}
% 
% Figure~\ref{fig:low} shows our OpenCL-specific rules that are used to apply OpenCL optimizations and to lower high-level concepts down to OpenCL-specific ones.
% Patterns that are known to the code generator are shown in bold in both Figure~\ref{fig:algo} and~\ref{fig:low}.
% %The code generation process is described separately in section~\ref{codegen}.
% 
% % \paragraph{Identity}
% % Figure~\ref{fig:low:id} shows the rewrite rule for the identity pattern.
% % We decided to only support code generation for the \pat{id-primitive} function
% % which is valid only on non-array data types.
% % % in our code generator which is the id function.
% % We explicitly use the \pat{map} pattern in order to implement the identity function for arrays.
% 
% \paragraph{Maps}
% The rule in Figure~\ref{fig:low:map} is used to produce OpenCL-specific map implementations that match the thread hierarchy of the OpenCL programming model.
% Our implementation maintains context information to ensure the thread hierarchy is respected.
% For instance, it is only legal to nest a \pat{map-local} inside a \pat{map-workgroup}.
% %Similarly, a \pat{map-global} pattern can not be nested inside a \pat{map-workgroup} pattern as it represents global threads not organized in OpenCL workgroups.
% 
% % The \emph{pre} and \emph{post} conditions reflect the hierarchy presented in Figure~\ref{fig:map} corresponding to the OpenCL programming model.
% % For instance for \pat{map-workgroup}, we have the pre condition that we must not be within a \pat{map-workgroup} ($\overline{mwg}$) already.
% 
% %Some rules can only be activated if certain conditions are true.
% %We use the syntax $[pre:post]$ to represent the pre and post conditions of a rule.
% %The pre condition $pre$ corresponds to the list of conditions that must be true for the rule to be applied.
% %The post condition $post$ is set for any function bound to the pattern.
% %The $\overline{\rule[-.3\baselineskip]{0pt}{1.5ex}\hspace{1.5ex}}$, $\wedge$, and $\vee$ corresponds to the logical \emph{not}, \emph{and}, and \emph{or} operators respectively.
% 
% \paragraph{Reduction}
% There is only one lowering rule for reduction (Figure~\ref{fig:low:red}), which expresses the fact that the only OpenCL implementation known to the code generator is a sequential reduction.
% Possible parallel implementations of the reduction pattern are defined at a higher level by composition of other algorithmic patterns.
% To the best of our knowledge, all other existing high performance compilers treat the reduction directly as an irreducible primitive operation.
% The power of our approach is that the code generator implementation only needs to know about the simple sequential reduction.
% As a result, it is possible to explore different implementation for the reduction by simply applying different rules.
% 
% \paragraph{Reorder}
% Figure~\ref{fig:low:stride} presents the rule that reorders elements of an array.
% In our current implementation, we support two types of reordering:
% no reordering, represented by the \pat{id} identify function, and \pat{reorder-stride} which reorders elements with a certain stride $s$.
% As described earlier, the major use case for the stride reorder is to enable coalesced memory accesses.
% %Note that other types of reordering functions could be implemented easily within this framework such as user-defined reorder function.
% 
% \paragraph{Local/Global}
% Figure~\ref{fig:low:mem} shows two rules that enable GPU local memory usage.
% They express the fact that the result of a \pat{map-local} can always be stored in local memory or back in global memory.
% This holds since a \pat{map-local} always exists within a \pat{map-workgroup} for which the local memory is defined.
% These rules allow us to determine how the data is mapped to the GPU memory hierarchy.
% 
% 
% 
% \paragraph{Vectorization}
% Finally, Figure~\ref{fig:algo:vect} shows the vectorization rule.
% Vectorization is achieved by using the \pat{asVector} and corresponding \pat{asScalar} which changes the element type of an array and adjust the length accordingly.
% This rule is only allowed to be applied once to a given \pat{map(f)} pattern.
% This constrain can easily be checked by looking at the function's type; if it is a vector type, the rule cannot be applied.
% Another set of rules, not shown here for space reason, are used to propagate the $vect^n$ function recursively within $f$.
% 
% %\subsection{Automatic Rewriting Strategy}
% %\label{sec:search}
% %
% %The rules presented in this section define a search space of possible implementations.
% %In order to find the best possible low-level expression for a given target device, we have developed a simple automatic search strategy based loosely on Bandit-based optimization~\cite{demesmay09bandit}.
% %Note that our current search strategy is just designed to prove that it is possible to find good implementations.
% %We envision replacing this exploration strategy in the future by using machine-learning techniques to avoid having to search the space.
% %However, this is orthogonal to the work presented in this paper.
% %
% %Our search strategy starts with the high-level expression and determines all the valid rules that can be applied at this stage.
% %We use a Monte-Carlo method for evaluating the potential impact of each rule by randomly walking down the search tree.
% %The rule that will lead to the best performance following the Monte-Carlo descent is chosen and applied to the expression.
% %This process is repeated until we reach a terminal expression.
% %Note that in addition to selecting the rules, we also search at the same time for the parameters controlling our patterns such as the vector size for the $\pat{vect}^n$ pattern.
% %Using this simple strategy, we found that less than a thousand expressions were evaluated to reach a solution in most cases.
% 
% 
% 
% \subsection{Summary}
% 
% The power of our approach lies in the composition of our rules that produce complex low-level expressions from simple high-level expressions.
% Looking back at our motivation example in Figure~\ref{fig:codeex}, we see how a simple algorithmic pattern such as \pat{map} can effectively be derived into a low-level expression by applying the rules.
% This expression matches various hardware concepts expressible with the OpenCL programming model such as mapping computation and data to the GPU thread and memory hierarchy and vectorization.
% Each single rule encodes a simple, easy to understand, provable fact.
% By composition of the rules we systematically derive low-level expressions which are semantically equivalent to the high-level expressions by construction.
% This results in a powerful mechanism to safely explore the space of possible implementations.
% % Using these rules and our automatic search technique we now have a powerful mechanism to explore the space of possible program implementations.
% \from{PPoPP end}
% 
% \from{PLDI begin}
% \paragraph{Rewrite Rules (PLDI)}
% This section introduces our set of rewrite rules that transform high-level expressions written using our algorithmic patterns into semantically equivalent expressions.
% One goal of our approach is to keep each rule as simple as possible and only express one fundamental concept at a time.
% For instance the vectorization rule, as we will see, is the only place where we express the vectorization concept.
% This is different from most prior approaches that would produce a special vectorized version of different algorithmic patterns such as map or reduce.
% The superiority of our approach lies in the power of composition;
% many rules can be applied successively to produce expressions that compose hardware concepts or optimizations and that are provably correct by construction.
% 
% Similarly to our patterns, we distinguish between algorithmic and lowering rules.
% Algorithmic rules produce derivations that represent the different algorithmic choices and are shown in Figure~\ref{fig:algo}.
% %The lowering rules map expressions to hardware patterns.
% Figure~\ref{fig:low} shows our OpenCL-specific rules which map expressions to OpenCL patterns.
% %are expressible with the OpenCL programming model.
% Once the expression is in its lowest form, it is possible to produce OpenCL code for each single pattern easily with our code generator as described in the previous section.
% 
% 
% %\paragraph{Syntax and Rule Derivation}
% %Some rules can only be activated if certain conditions are true.
% %We use the syntax $[pre:post]$ to represent the pre and post conditions of a rule.
% %The pre condition $pre$ corresponds to the list of conditions that must be true for the rule to be applied.
% %The post condition $post$ is set for any function bound to the pattern.
% %The $\overline{\rule[-.3\baselineskip]{0pt}{1.5ex}\hspace{1.5ex}}$, $\wedge$, and $\vee$ corresponds to the logical \emph{not}, \emph{and}, and \emph{or} operators respectively.
% 
% %We use leftmost derivation when applying the rules, which means that the leftmost non-terminal is always derived first.
% 
% 
% 
% 
% 
% \begin{figure}[t]
% \centering
% 
% % \begin{subfigure}[b]{1\linewidth}
% % \begin{mdframed}
% % $$
% % \begin{array}{llll}
% %  \langle primitive \rangle & \pat{id} & \rightarrow & \pat{\textbf{id-primitive}}\\
% %  \langle T[\ ] \rangle     & \pat{id} & \rightarrow & \pat{map(id)}\\  
% % \end{array}
% % $$
% % \end{mdframed}
% %   \caption{Identity}
% %   \label{fig:low:id}
% % \end{subfigure}
% 
% \begin{subfigure}[b]{1\linewidth}
% \begin{mdframed}
% \vspace{-.5em}
% $$
% \begin{array}{lllll}
% \pat{map(f)} & \rightarrow & \pat{\textbf{map-workgroup}(f)} & | & \pat{\textbf{map-local}(f)}\\
%  & | & \pat{\textbf{map-global}(f)} & | & \pat{\textbf{map-seq}(f)}\\          
% \end{array}
% $$
% \end{mdframed}
%   \caption{Map}
%   \label{fig:low:map}
% \end{subfigure}
% 
% \vspace{\ruleSpace}
% \begin{subfigure}[b]{1\linewidth}
% \begin{mdframed}
% $$
% \begin{array}{lll}
%   \pat{reduce(f,z)} & \rightarrow & \pat{\textbf{reduce-seq}(f,z)}% & [\hspace{.2em}mlc\vee mln:\ \hspace{.2em}]\\   
% \end{array}
% $$
% \end{mdframed}
%   \caption{Reduction}
%   \label{fig:low:red}
% \end{subfigure}
% 
% \vspace{\ruleSpace}
% \begin{subfigure}[b]{1\linewidth}
% \begin{mdframed}
% $$
% \begin{array}{lllll}
%   \pat{reorder}  & \rightarrow & \pat{\textbf{reorder-stride}}^s & | & \pat{id}
% \end{array}
% $$
% \end{mdframed}
%   \caption{Stride accesses or normal accesses}
%   \label{fig:low:stride}
% \end{subfigure}
% 
% \vspace{\ruleSpace}
% \begin{subfigure}[b]{1\linewidth}
% \begin{mdframed}
% \vspace{.2em}
% $$
% \begin{array}{lrl}
%   \pat{\textbf{map-local}(f)} & \rightarrow & \pat{\textbf{toGlobal}(\textbf{map-local}(f))}\\  
%   \pat{\textbf{map-local}(f)} & \rightarrow & \pat{\textbf{toLocal}(\textbf{map-local}(f))}\\  
% \end{array}
% $$
% \end{mdframed}
%   \caption{Local/Global memory}
%   \label{fig:low:mem}
% \end{subfigure}
% 
% \vspace{\ruleSpace}
% \begin{subfigure}[b]{1\linewidth}
% \begin{mdframed}
% \vspace{-.75em}
% $$
% \begin{array}{lrll}
% \pat{map(f)} & \rightarrow & \pat{\textbf{asScalar}} \circ \pat{map(\textbf{vect}}^n\pat{(f))} \circ \pat{\textbf{asVector}}^n
% \end{array}
% $$
% \end{mdframed}
%   \caption{Vectorization}
%    \label{fig:algo:vect}
% \end{subfigure}
% 
% \vspace{-2em}
% \caption{OpenCL-specific rules. Bold patterns are known to the code generator.}
% \label{fig:low}
% \end{figure}
% 
% 
% 
% 
% \subsection{Algorithmic Rules}
% 
% \newcommand{\Reduce}{\text{\textit{reduce}}\xspace}
% \newcommand{\PartRed}{\text{\textit{part-red}}\xspace}
% \newcommand{\RedSeq}{\text{\textit{reduce-seq}}\xspace}
% \newcommand{\Map}{\text{\textit{map}}\xspace}
% \newcommand{\MapSeq}{\text{\textit{map-seq}}\xspace}
% \newcommand{\MyJoin}{\text{\textit{join}}\xspace}
% \newcommand{\MySplit}[1]{\text{\textit{split}}^{#1}\xspace}
% 
% \begin{figure*}[t]
% \begin{align}
%   \text{\textit{asum}}(\vec{x})
%   & =\hspace{.2em} \Reduce(+, 0) \circ \Map(abs, \vec{x})
%    \overset{\ref{fig:algo:red}}{\hspace{.2em}=\hspace{.2em}}
%       \Reduce(+, 0) \circ \MyJoin \circ \Map(\PartRed(+, 0)) \circ \MySplit{n} \circ \Map(abs, \vec{x})\\[-.5em]
%   & \overset{\ref{fig:algo:splitjoin}}{=\hspace{.2em}}
%       \Reduce(+, 0) \circ \MyJoin \circ \Map(\PartRed(+, 0)) \circ \MySplit{n} \circ \MyJoin \circ \Map(\Map(abs)) \circ \MySplit{n}(\vec{x})\\[-.5em]
%   & \overset{\ref{fig:algo:simpl}}{=\hspace{.2em}}
%       \Reduce(+, 0) \circ \MyJoin \circ \Map(\PartRed(+, 0)) \circ \Map(\Map(abs)) \circ \MySplit{n}(\vec{x})\\[-.5em]
%   & \overset{\ref{fig:algo:fusion}}{=\hspace{.2em}}
%       \Reduce(+, 0) \circ \MyJoin \circ \Map(\PartRed(+, 0) \circ \Map(abs)) \circ \MySplit{n}(\vec{x})\\[-.5em]
%   & \overset{\ref{fig:low:map}}{=\hspace{.2em}}
%       \Reduce(+, 0) \circ \MyJoin \circ \Map(\PartRed(+, 0) \circ \MapSeq(abs)) \circ \MySplit{n}(\vec{x})\\[-.5em]
%   & \overset{\hspace{-1.3em}\ref{fig:algo:red}\& \ref{fig:low:red}\hspace{-1.1em}}{=\hspace{.2em}}
%       \Reduce(+, 0) \circ \MyJoin \circ \Map(\RedSeq(+, 0) \circ \MapSeq(abs)) \circ \MySplit{n}(\vec{x})\\[-.5em]
%   & \overset{\ref{fig:algo:fusion}}{=\hspace{.2em}}
%       \Reduce(+, 0) \circ \MyJoin \circ \Map(\RedSeq(\lambda\ acc, a: acc + abs(a), 0) \circ \MySplit{n}(\vec{x})
% \end{align}
% \vspace{-2em}
% \caption{Derivation for \emph{asum}$(\vec{x})$ to a fused parallel version.
%   The numbers %above the equality sign
%   refer to the rules from Figure~\ref{fig:algo} and Figure~\ref{fig:low}.
% }
% \label{fig:derivation}
% \end{figure*}
% 
% % The first three rules (see Figure~\ref{fig:algo:identity},~\ref{fig:algo:iterate} and~\ref{fig:algo:reorder}) express the properties of the identity, iterate and reorder pattern.
% 
% % \paragraph{Identity}
% % The identity rule in Figure~\ref{fig:algo:identity} specifies that it is always valid to compose any function with the identity function \pat{id}.
% % The \pat{id} function can acts as a copy operation; this is useful to express copy to local memory for instance when composed with the corresponding lowering rule.
%  
% \paragraph{Iterate decomposition}
% The rule~\ref{fig:algo:iterate} expresses the fact that an iteration can be decomposed into several iterations.
% 
% \paragraph{Reorder commutativity}
% Figure~\ref{fig:algo:reorder} shows that if the data can be reordered arbitrarily it does not matter if we apply a function $f$ to each element before or after the reordering.
% 
% \paragraph{Split-join}
% The split-join rule in Figure~\ref{fig:algo:splitjoin} partitions a map into two maps.
% This allows us to nest map patterns in each other and, thus, \emph{maps} the computation to the thread hierarchy of the OpenCL programming model.
% % such as \pat{map-workgroup(map-local(f))} as seen in our motivation example (Figure~\ref{fig:codeex}).
% 
% 
% \paragraph{Reduction}
% The reduction (and associated partial reduction) in Figure~\ref{fig:algo:red} is currently our most complex rule but also the most powerful one.
% It expresses the reduction function as a composition of other primitive functions, which is a fundamental aspect of our work.
% From the algorithmic point of view we first define a partial reduction pattern \pat{part-red}.
% This partial reduction reduces an array of $n$ elements to an array of $m$ elements where $1 \leq m < n$.
% The reduction can be derived in a partial reduction combined with a full reduction which ensures we end up with one unique element.
% % Another possible derivation consists in iterating a partial reduction until a full reduction is achieved (this is what $\infty$ represents).
% %Note that our definition of \pat{reduce} remains correct since the result of a partial reduction is always composed with the reduction to ensure we end up with one unique element.
% 
% \paragraph{Partial Reduction}
% The first possible derivation for partial reduction, in Figure~\ref{fig:algo:red}, leads to the full reduction ($m=1$).
% %The first two possible derivations for partial reduction, in figure~\ref{fig:algo:red}, lead to the $id$ function or the full reduction which means $m=n$ or $m=1$ respectively.
% The next possible derivation expresses the fact that it is possible to reorder the elements to be reduced, expressing the commutativity property of our definition of reduction.
% The third derivation is actually the only place where parallelism is expressed in the definition of our reduction pattern.
% This rule expressed the fact that it is valid to partition the input elements first and then reduce them independently.
% Finally, the last possible derivation expresses the notion that it is possible to perform a partial reduction with an iterative process by repetitively applying the same partial reduction function.
% This concept is very important when considering how the reduction function is typically implemented on a GPU (iteratively reducing within a workgroup using the local memory).
% 
% 
% \paragraph{Simplification Rules}
% Figure~\ref{fig:algo:simpl} shows our simplification rules.
% They express the fact that consecutive \pat{split}-\pat{join} pairs and \pat{asVector}-\pat{asScalar} pairs are equivalent to the identity. % function \emph{id}.
% 
% \paragraph{Fusion Rules}
% Finally, our fusion rules are shown in Figure~\ref{fig:algo:fusion}.
% The first rule fuses the functions applied by two consecutive maps.
% The second rule fuses the map-reduce pattern by creating a lambda function that is the results of merging function $f$ and $g$ from the original reduction and map respectively.
% This rule only applies to the sequential version since this is the only implementation not requiring the associativity property required by the more generic $\pat{reduce}$ pattern.
% When generating code, these rules in effect allow us to fuse the implementation of the different functions and avoid having to store temporary results.
% % More sophisticated rules exist in the functional programming community~\cite{fusion} which we want to incorporate in the future.
% % The functional community has also studied rules for fusion (a.k.a. \emph{deforestation}) from a theoretical~\cite{} as well as a more practical point of view~\cite{jones01playing}.
% The functional programming community has studied more generic rules for fusion~\cite{coutts07streamfusion,jones01playing}.
% %More generic rules for fusion have been studied by the functional programming community~\cite{coutts07streamfusion,jones01playing}.
% However, as we currently focus on a restricted set of patterns our simpler fusion rules have, so far, proven to be sufficient.
% 
% 
% 
% 
% 
% 
% \subsection{OpenCL-Specific Rules}
% 
% Figure~\ref{fig:low} shows our OpenCL-specific rules that are used to apply OpenCL optimizations and to lower high-level concepts down to OpenCL-specific ones.
% Patterns that are known to the code generator are shown in bold.% in both Figure~\ref{fig:algo} and~\ref{fig:low}.
% %The code generation process is described separately in section~\ref{codegen}.
% 
% % \paragraph{Identity}
% % Figure~\ref{fig:low:id} shows the rewrite rule for the identity pattern.
% % We decided to only support code generation for the \pat{id-primitive} function
% % which is valid only on non-array data types.
% % % in our code generator which is the id function.
% % We explicitly use the \pat{map} pattern in order to implement the identity function for arrays.
% 
% \paragraph{Maps}
% The rule in Figure~\ref{fig:low:map} is used to produce OpenCL-specific map implementations that match the thread hierarchy of OpenCL.
% Our implementation maintains context information to ensure the hierarchy is respected.
% For instance, it is only legal to nest a \pat{map-local} inside a \pat{map-workgroup}.
% %Similarly, a \pat{map-global} pattern can not be nested inside a \pat{map-workgroup} pattern as it represents global threads not organized in OpenCL workgroups.
% 
% % The \emph{pre} and \emph{post} conditions reflect the hierarchy presented in Figure~\ref{fig:map} corresponding to the OpenCL programming model.
% % For instance for \pat{map-workgroup}, we have the pre condition that we must not be within a \pat{map-workgroup} ($\overline{mwg}$) already.
% 
% %Some rules can only be activated if certain conditions are true.
% %We use the syntax $[pre:post]$ to represent the pre and post conditions of a rule.
% %The pre condition $pre$ corresponds to the list of conditions that must be true for the rule to be applied.
% %The post condition $post$ is set for any function bound to the pattern.
% %The $\overline{\rule[-.3\baselineskip]{0pt}{1.5ex}\hspace{1.5ex}}$, $\wedge$, and $\vee$ corresponds to the logical \emph{not}, \emph{and}, and \emph{or} operators respectively.
% 
% \paragraph{Reduction}
% There is only one lowering rule for reduction (Figure~\ref{fig:low:red}), which expresses the fact that the only implementation known to the code generator is a sequential reduction.
% Parallel implementations are defined at a higher level by composition of other algorithmic patterns.
% Most existing high performance compilers treat the reduction directly as an irreducible primitive operation.
% %The power of our approach is that the code generator implementation only needs to know about the simple sequential reduction.
% %As a result, it is possible to explore different implementation for the reduction by simply applying different rules.
% With our approach it is possible to explore different implementations for the reduction by simply applying different rules.
% 
% \paragraph{Reorder}
% Figure~\ref{fig:low:stride} presents the rule that reorders elements of an array.
% In our current implementation, we support two types of reordering:
% no reordering, represented by the \pat{id} function, and \pat{reorder-stride}, which reorders elements with a certain stride $s$.
% As described earlier, the major use case for the stride reorder is to enable coalesced memory accesses.
% %Note that other types of reordering functions could be implemented easily within this framework such as user-defined reorder function.
% 
% \paragraph{Local/Global}
% Figure~\ref{fig:low:mem} shows two rules that enable GPU local memory usage.
% They express the fact that the result of a \pat{map-local} can always be stored in local memory or back in global memory.
% This holds since a \pat{map-local} always exists within a \pat{map-workgroup} for which the local memory is defined.
% These rules allow us to determine how the data is mapped to the GPU memory hierarchy.
% 
% 
% 
% \paragraph{Vectorization}
% Figure~\ref{fig:algo:vect} shows the vectorization rule.
% Vectorization is achieved by using the \pat{asVector} and corresponding \pat{asScalar}, which changes the element type of an array and adjust the length accordingly.
% This rule is only allowed to be applied once to a given \pat{map(f)} pattern.
% This constrain can easily be checked by looking at the function's type. %; if it is a vector type, the rule cannot be applied.
% Another set of rules, not shown here for space reason, is used to propagate the $vect^n$ function recursively within $f$.
% 
% %\subsection{Automatic Rewriting Strategy}
% %\label{sec:search}
% %
% %The rules presented in this section define a search space of possible implementations.
% %In order to find the best possible low-level expression for a given target device, we have developed a simple automatic search strategy based loosely on Bandit-based optimization~\cite{demesmay09bandit}.
% %Note that our current search strategy is just designed to prove that it is possible to find good implementations.
% %We envision replacing this exploration strategy in the future by using machine-learning techniques to avoid having to search the space.
% %However, this is orthogonal to the work presented in this paper.
% %
% %Our search strategy starts with the high-level expression and determines all the valid rules that can be applied at this stage.
% %We use a Monte-Carlo method for evaluating the potential impact of each rule by randomly walking down the search tree.
% %The rule that will lead to the best performance following the Monte-Carlo descent is chosen and applied to the expression.
% %This process is repeated until we reach a terminal expression.
% %Note that in addition to selecting the rules, we also search at the same time for the parameters controlling our patterns such as the vector size for the $\pat{vect}^n$ pattern.
% %Using this simple strategy, we found that less than a thousand expressions were evaluated to reach a solution in most cases.
% 
% \begin{figure*}[t]
% \captionsetup[subfigure]{justification=justified,singlelinecheck=false}
% 
% \begin{subfigure}[b]{\linewidth}
% \vspace{.4em}
% \begin{minipage}{.02\linewidth}
% \caption{\hspace{-.2em}\parbox[t]{.7\linewidth}{\centering Nvidia\\ GPU}}
% \label{fig:expr:autoNv}
% \end{minipage}
% \hfill
% \begin{minipage}{.93\linewidth}
% \begin{lstlisting}[mathescape]
% def asum($\vec{x}$) = reduce-seq o join o join o map-workgroup(
%   toGlobal(map-local(iterate-7(reduce-seq(plus, 0)) o reduce-seq(absAndPlus, 0))) o reorder-stride
%  ) o split-128 o split-2048($\vec{x}$)
% \end{lstlisting}
% % \begin{lstlisting}[mathescape]
% % def asum($\vec{x}$) = reduce-seq o join o join o map-workgroup(
% %   toGlobal(map-local(iterate-7(map-seq(id) o reduce-seq(plus, 0)) o reduce-seq(absAndPlus, 0))) o reorder-stride
% %  ) o split-128 o split-2048($\vec{x}$)
% % \end{lstlisting}
% \end{minipage}
% \end{subfigure}
% 
% \begin{subfigure}[b]{\linewidth}
% \vspace{0em}
% \begin{minipage}{.02\linewidth}
% \caption{\hspace{-.2em}\parbox[t]{.7\linewidth}{\centering AMD\\ GPU}}
% \label{fig:expr:autoAmd}
% \end{minipage}
% \hfill
% \begin{minipage}{.93\linewidth}
% \begin{lstlisting}[mathescape]
% def asum($\vec{x}$) = join o asScalar o join o map-workgroup(
%    map-local( map-seq(vect-2(id)) o reduce-seq(vect-2(absAndPlus), vect-2(0))) o reorder-stride
%  ) o split-128 o asVector-2 o split-4096($\vec{x}$)
% \end{lstlisting}
% \end{minipage}
% \end{subfigure}
% 
% \begin{subfigure}[b]{\linewidth}
% \vspace{0em}
% \begin{minipage}{.02\linewidth}
% \caption{\hspace{-.2em}\parbox[t]{.7\linewidth}{\centering Intel\\ CPU}}
% \label{fig:expr:autoCPU}
% \end{minipage}
% \hfill
% \begin{minipage}{.93\linewidth}
% \begin{lstlisting}[mathescape]
% def asum($\vec{x}$) = join o map-workgroup( join o asScalar o map-local(
%    map-seq(vect-4(id)) o reduce-seq(vect-4(absAndPlus), vect-4(0))
%  ) o asVector-4 o split-32768 ) o split-32768($\vec{x}$)
% \end{lstlisting}
% \end{minipage}
% \end{subfigure}
% 
% \caption{Low-level expressions performing the sum of absolute values.
%          These expressions are automatically derived by our system from the high-level expression \textit{asum}$(\vec{x}) = $\textit{ reduce}(+, 0) $\circ$ \textit{map}(\textit{abs}, $\vec{x}$).
%        }
% \label{fig:expr}
% \end{figure*}
% 
% 
% \subsection{Example: Deriving a Fused Implementation}
% \label{sec:example}
% 
% 
% 
% 
% 
% To achieve good performance it is in general beneficial to avoid storing intermediate results.
% Rule~\ref{fig:algo:fusion} allows us to apply this principle and fuse two patterns into one, thus, avoiding intermediate results.
% Figure~\ref{fig:derivation} shows how we can derive a fused version for calculating the sum of absolute value, \emph{asum}, from the high-level expression written by the programmer.
% We write the derivation as a sequence of equations using a slightly more mathematical notation.
% The numbers above the equality sign refer to the rules applied.
% 
% We start by applying the reduction rule~\ref{fig:algo:red} twice:
% first to replace \pat{reduce} with \pat{reduce}~$\circ$~\pat{part-red} and then a second time to expand \pat{par-red}.
% To get (2) we expand \pat{map}, which can be simplified by removing the two corresponding \pat{join} and \pat{split} patterns.
% In the step from (3) to (4) two \pat{map} patterns are fused and in the next step the nested \pat{map} is lowered into the \pat{map-seq} pattern to obtain (5).
% By first transforming \pat{part-red} back into \pat{reduce} (using rule~\ref{fig:algo:red}) and then applying the lowering rule~\ref{fig:low:red} we get (6).
% Finally, we apply rule~\ref{fig:algo:fusion} to fuse the \pat{map-seq} and \pat{reduce-seq} into a single \pat{reduce-seq}.
% This sequence of transformations results in expression (7), which allows for a more optimal implementation since no temporary storage is required for the intermediate result.
% 
% 
% 
% 
% 
% 
% 
% \subsection{Summary}
% 
% The power of our approach lies in the composition of our rules that produce complex low-level expressions from simple high-level expressions.
% Looking back at our example in Figure~\ref{fig:codeex}, we see how a simple algorithmic pattern can effectively be derived into a low-level expression by applying the rules.
% This expression matches hardware concepts expressible with OpenCL such as mapping computation and data to the thread and memory hierarchy. % and vectorization.
% Each single rule encodes a simple, easy to understand, provable fact.
% By composition of the rules we systematically derive low-level expressions which are semantically equivalent to the high-level expressions by construction.
% This results in a powerful mechanism to safely explore the space of possible implementations.
% % Using these rules and our automatic search technique we now have a powerful mechanism to explore the space of possible program implementations.
% \from{PLDI end}




\section{Code Generator \& Implementation Details}
In this section we discuss how an low-level expression comprising of patterns from \autoref{section:patterns} and possibly derived using the rewrite rules from \autoref{section:rules} is turned into \OpenCL code.
We will see, that this process is surprisingly simple and straightforward.
This is due to the fact, that all complex decisions regarding optimizations are made at an earlier stage: when applying the rewrite rules.
This design was chosen deliberately and follows the principle of separation of concerns and keeps the implementation of the code generator simple.
The expression used as input for the code generator explicitly specifies every important detail of the generated \OpenCL implementation, so that for every expression there is a one-to-one mapping to \OpenCL code.

We will start our discussion by looking back at the patterns defined in \autoref{section:patterns} and how \OpenCL code is generated for them.
We will see, that there are patterns for which it is not possible to generate \OpenCL code.
These expressions do not specify the \OpenCL implementation detailed enough.
We can use the rewrite rules presented in \autoref{section:rules} to transform the expression until, finally, the expression is precise enough for the code generator.

We then shift our focus to the implementation of the type system and how this helps to implement a static memory allocator.

Finally, we will give some details specifying which software infrastructure we used in our implementation.

\subsection{Generating \OpenCL Code for Patterns}
When generating \OpenCL code from an expression composed of patterns the code generator traverses the expression.
It follows the data flow and emits code for each pattern it visits.

\begin{table}[t]
\centering
\begin{tabular}{llll}
\toprule
    \multicolumn{2}{c}{\tabhead{Algorithmic Patterns}}
  & \multicolumn{2}{c}{\tabhead{\OpenCL Patterns}}\\
\midrule
 \map &
  \textbf{\zip} &
    \textbf{\mapWorkgroup} &
      \textbf{\reduceSeq}\\
 \reduce&
  \textbf{\splitN}&
    \textbf{\mapLocal}&
      \textbf{\toLocal}\\
 \reorder&
  \textbf{\join} &
    \textbf{\mapGlobal}&
      \textbf{\toGlobal}\\
 &
  \textbf{\iterateN} &
    \textbf{\mapWarp}&
      \textbf{\reorderStride}\\
 & &
    \textbf{\mapLane} &
      \textbf{\asVector}\\
 & & \textbf{\mapSeq} &
        \textbf{\asScalar}\\
 & & & \textbf{\vect}\\
\bottomrule
\end{tabular}
\caption{Overview of all algorithmic and \OpenCL patterns.}
\label{fig:patterns:generation}
\end{table}

\autoref{fig:patterns:generation} shows all patterns introduced in \autoref{section:patterns}.
The code generator does not know how to generate \OpenCL code for all patterns, for example are there many different options to implement the \reduce pattern in \OpenCL, as we discussed in \autoref{section:reduce:case-study}.
Therefore, the code generator would have to make a choice which implementation to pick.
We want to avoid such situations, as these complicates the implementation of the code generator and limits its flexibility and the performance of the generated code.
In our approach this decision about the implementation of \reduce has to be made before the code generator is invoked by applying the rewrite rules presented in \autoref{section:rules} which allow to safely derive specialized low-level implementations from high-level expressions.

The code generator generates code only for the patterns highlighted in bold in \autoref{fig:patterns:generation} (all patterns in the last three columns).
The three patterns in the first column: \map, \reduce, and \reorder have to be eliminated from an expression before the code generation process can begin.

We will now discuss the code generation process for all highlighted patterns from \autoref{fig:patterns:generation} in more detail.

\paragraph{Zip}
The code generator emits no \OpenCL code for the \zip pattern.
Instead \zip's type has an effect how code for following patterns is generated.
Let us look at the following example, where the \zip and \mapGlobal patterns are used together:
\begin{align}
  \map\ (+)\ (\zip\ \vec{x}\ \vec{y})
\end{align}
When processing this expression the \zip pattern is visited first.
The type of \zip makes it explicit to the code generator that, when emitting code for the following \mapGlobal, two elements -- one element from each array -- have to be read together.
In the implementation the code generator will investigate the type of the input array before emitting code for the \mapGlobal.

\paragraph{Split and Join}
Similar to the \zip pattern no \OpenCL code is emitted for neither \splitN nor \join.
By encoding the size of arrays in the type system the information how the data was shaped by the \splitN and \join patterns is passed to following patterns.
This information is used later when generating the \OpenCL code for performing \OpenCL memory accesses. 
We will discuss the type system implementation in more detail in \autoref{section:typeSystem}.

\paragraph{Iterate}
\autoref{lst:iterate:impl} shows the structure of the \OpenCL code generated for the \iterateN pattern.
%
\begin{lstlisting}[%                                                             
caption={Structure of the \OpenCL code emitted for the \iterateN pattern.},%
numbers=left,%
float=tb,
label={lst:iterate:impl}]
int size = 128;
local float* in = array0;
local float* out = ((7 & 1) != 0) ? array1 : array2;
for (int i = 0; i < 7; i +=1) {$\label{lst:iterate:impl:for}$
  ... // code emitted for nested pattern$\label{lst:iterate:nested}$
  in  = (out == array2) ? array2 : array1;$\label{lst:iterate:impl:in}$
  out = (out == array2) ? array2 : array1;$label{lst:iterate:impl:out}$
  size = size / 2;$\label{lst:iterate:impl:size}$
}
\end{lstlisting}
%
This code segment is actually taken from the parallel reduction example discussed in \autoref{eq:reduce11}.
A for loop is generated (line~\ref{lst:iterate:impl:for}) for performing the actual iterations.
Two pointers (\code{in} and \code{out}) are swapped after each iteration (lines~\ref{lst:iterate:impl:in} and~\ref{lst:iterate:impl:out}) and used in the code generated for the nested pattern which is emitted inside the body of the for loop (line~\ref{lst:iterate:nested}).
Because this code segment is taken from the parallel reduction example, the \code{size} variable, which represents the size of the current input array, is reduced by half after each iteration step (line~\ref{lst:iterate:impl:size}).

\paragraph{Parallel OpenCL Maps}
The generation of \OpenCL code for one of the \map patterns: \mapWorkgroup, \mapLocal, \mapGlobal, \mapWarp, and \mapLane is rather straightforward.

\autoref{lst:mapWG:impl} shows the structure of the \OpenCL code generated for the \mapWorkgroup pattern.
%
\begin{lstlisting}[%                                                             
caption={Structure of the \OpenCL code emitted for the \mapWorkgroup pattern.},%
numbers=left,%
float=tb,
label={lst:mapWG:impl}]
for (int wg_id = get_group_id(0); wg_id < size;
     wg_id += get_num_groups(0)) {
  ... // code emitted for nested pattern
}
return;
\end{lstlisting}
%
A for loop is emitted where the loop variable (\code{wg\_id}) represents the work-group id in \OpenCL.
The loop variable is used by the nested pattern inside the loop to ensure that the correct index is accessed in memory.
After the for loop an synchronization mechanism is emitted, in this case \code{return}.
As global synchronization between work-items is not possible in a single kernel in \OpenCL, the only way for global synchronization is to terminate the \OpenCL kernel and continue the processing in a new one. 

\autoref{lst:mapLocal:impl} shows the \OpenCL code generated for the \mapLocal pattern.
%
\begin{lstlisting}[%                                                             
caption={Structure of the \OpenCL code emitted for the \mapLocal pattern.},%
numbers=left,%
float=tb,
label={lst:mapLocal:impl}]
for (int l_id_id = get_local_id(0); l_id < size;
     l_id += get_local_size(0)) {
  ... // code emitted for nested pattern
}
barrier(CLK_LOCAL_MEM_FENCE);
\end{lstlisting}
%
Here the \OpenCL \code{get\_local\_id(0)} function is used to obtain the id of the work-item inside the work-group.
The \code{barrier} function is used for synchronization.

The \OpenCL code emitted for \mapGlobal, \mapWarp, and \mapLane is similar reflecting the specific way of obtaining the loop variable and synchronization mechanism in \OpenCL.
As there is no synchronization required between work-item of a single warp, the \mapLane pattern emits no synchronization statement.

\paragraph{Sequential Map and Reduction}
The \OpenCL implementations of the sequential \mapSeq and \reduceSeq patterns are shown in \autoref{lst:mapSeq:impl} and \autoref{lst:redSeq:impl}.
%
\begin{lstlisting}[%                                                             
caption={Structure of the \OpenCL code emitted for the \mapSeq pattern.},%
numbers=left,%
float=tb,
label={lst:mapSeq:impl}]
for (int i = 0; i < size; ++i) {
  output[out_index] = f(input[in_index]);
}
\end{lstlisting}
%
%
\begin{lstlisting}[%                                                             
caption={Structure of the \OpenCL code emitted for the \reduceSeq pattern.},%
numbers=left,%
float=tb,
label={lst:redSeq:impl}]
float acc = 0.0f;
for (int i = 0; i < size; ++i) {
  acc = f(acc, input[in_index]);
}
output[out_index] = acc;
\end{lstlisting}
%
Both implementations are straightforward.
The \mapSeq implementation applies its customizing function to each element of the input array and stores the outputs in an output array.
The \reduceSeq implementation uses an accumulation variable to accumulate the result of the reduction while it iterates through the input array.
After the for loop the result is stored in the output array.

The most complicated bit is the generation of the input and output indices: \code{in\_index} and \code{out\_index}.
To understand how this works let us investigate an example which combines the previous code segments.
\autoref{lst:mapTogether:impl} shows the code generated for the expression:
\begin{align*}
  &\join \circ \mapWorkgroup\ \big(\\
  &\qquad\join \circ \mapLocal\ (\mapSeq\ id) \circ \splitN\ 1\big) \circ \splitN\ 128
\end{align*}
%
\begin{lstlisting}[%                                                             
caption={Structure of the \OpenCL code emitted for the \mapLocal pattern.},%
numbers=left,%
float=tb,
label={lst:mapTogether:impl}]
for (int wg_id = get_group_id(0); wg_id < size / 128;
     wg_id += get_num_groups(0)) {
  for (int l_id_id = get_local_id(0); l_id < 128;
       l_id += get_local_size(0)) {
    for (int i = 0; i < 1; ++i) {
      output[wg_id * 128 + l_id + i]$\label{lst:mapTogether:impl:output}$
        = id(input[wg_id * 128 + l_id + i]);$\label{lst:mapTogether:impl:input}$
    }
  }
  barrier(CLK_LOCAL_MEM_FENCE);
}
return;
\end{lstlisting}
%
The input and output indices (line~\ref{lst:mapTogether:impl:output} and~\ref{lst:mapTogether:impl:input})

\begin{lstlisting}[%                                                             
caption={Structure of the \OpenCL code emitted for the \mapLocal pattern.},%
numbers=left,%
float=tb,
label={lst:mapTogether:impl2}]
for (int wg_id = get_group_id(0); wg_id < size / 128;
     wg_id += get_num_groups(0)) {
  int l_id = get_local_id(0);
  output[wg_id * 128 + l_id + 0]
    = id(input[wg_id * 128 + l_id + 0]);
  barrier(CLK_LOCAL_MEM_FENCE);
}
return;
\end{lstlisting}


\paragraph{{\footnotesize to}Local and {\footnotesize to}Global}

\paragraph{Reorder-Stride}

\paragraph{{\footnotesize as}Vector, {\footnotesize as}Scalar, and Vectorize}



\subsection{The Type System and Static Memory Allocation}
\label{section:typeSystem}

\subsection{Implementation Details}
Our system is implemented in \Cpp, using the template system and support for lambda functions. 
When generating code for a given low-level expression, two basic steps are performed.
First, we use the Clang/LLVM compiler infrastructure to parse the expression and produce an abstract syntax tree for it.
Second, we traverse the tree and emit code for every function call representing one of our low-level hardware patterns.

As part of the first step, we have developed a type system which plays a dual role.
Firstly, it prevents the user, or a rewrite rule, to produce an expression that is not correct.
Secondly, the type system encodes informations that are necessary for code generation, such as memory address space and array size information, which are used to allocate memory.

The design of our code generator is straightforward since no optimization decisions are made at this stage.
We avoid performing complex analysis of the code which makes our design very different compared to traditional optimizing compilers.



\section{Summary}
\todo{List contributions and benefits of using the pattern code generator approach}

